%% ========================================================================

%% Zhu Shijie: Suanxue Qimeng (算學啓蒙) — Bilingual Edition

%% English-Chinese Parallel Text · Volume I (卷上) with General Principles (總括)

%% ========================================================================

\documentclass[10pt,a4paper]{article}



%% -------- Core Packages --------

\usepackage{fontspec}

\usepackage{polyglossia}

\usepackage{amsmath,amssymb}

\usepackage{geometry}

\usepackage{graphicx}

\usepackage{xcolor}

\usepackage{titlesec}

\usepackage{fancyhdr}

\usepackage{enumitem}

\usepackage{array,booktabs,longtable}

\usepackage{hyperref}

\usepackage{xeCJK}

\usepackage{etoolbox}



%% -------- Page Geometry --------

\geometry{

  paperwidth=210mm,paperheight=297mm,

  left=12mm,right=12mm,top=25mm,bottom=25mm,

  headheight=14pt,footskip=20pt

}



%% -------- Language Setup --------

\setmainlanguage{english}



%% -------- Font Configuration --------

\setmainfont{Times New Roman}

\setsansfont{Arial}

\setmonofont{Courier New}



%% CJK Support via xeCJK

\setCJKmainfont{Microsoft YaHei}

\setCJKsansfont{Microsoft YaHei}

\setCJKmonofont{Microsoft YaHei}



%% -------- Colors --------

\definecolor{titlecolor}{RGB}{45,55,72}

\definecolor{sectioncolor}{RGB}{55,65,81}

\definecolor{notecolor}{RGB}{120,53,15}

\definecolor{verselinecolor}{RGB}{80,60,40}

\definecolor{rulecolor}{RGB}{180,170,150}



%% -------- Section Formatting --------

\titleformat{\section}

  {\normalfont\large\bfseries\color{sectioncolor}}

  {\thesection}{0.5em}{}

\titleformat{\subsection}

  {\normalfont\normalsize\bfseries\color{sectioncolor}}

  {\thesubsection}{0.5em}{}



\titlespacing*{\section}{0pt}{1.5ex plus 0.5ex minus 0.2ex}{1ex plus 0.2ex}

\titlespacing*{\subsection}{0pt}{1ex plus 0.3ex minus 0.1ex}{0.5ex plus 0.1ex}



%% -------- Header/Footer --------

\pagestyle{fancy}

\fancyhf{}

\fancyhead[L]{\small\textit{Suanxue Qimeng · Bilingual Edition}}

\fancyhead[R]{\small\thepage}

\renewcommand{\headrulewidth}{0.3pt}

\renewcommand{\footrulewidth}{0pt}



%% -------- Custom Commands --------

\newcommand{\longline}{\noindent\rule{\textwidth}{0.3pt}}

\newcommand{\shortline}{\noindent\rule{0.5\textwidth}{0.3pt}\par}

\newcommand{\cnote}[1]{\textcolor{notecolor}{\small\textit{#1}}}



%% Bilingual parallel column commands

\newlength{\cnwidth}

\newlength{\enwidth}

\setlength{\cnwidth}{0.47\textwidth}

\setlength{\enwidth}{0.47\textwidth}



\newcommand{\bilingual}[2]{%

  \noindent%

  \begin{minipage}[t]{\cnwidth}%

    #1%

  \end{minipage}%

  \hfill\vrule width 0.4pt\hfill%

  \begin{minipage}[t]{\enwidth}%

    #2%

  \end{minipage}%

  \par\smallskip%

}



\newcommand{\bilingualsection}[2]{%

  \section*{#1 · #2}

  \addcontentsline{toc}{section}{#1 · #2}

}



\newcommand{\bilingualsubsection}[2]{%

  \subsection*{#1 · #2}

}



%% Verse-style rhyme environment

\newenvironment{rhymecn}{%

  \begin{quote}\itshape\color{verselinecolor}

}{\end{quote}}



%% Explanatory note

\newcommand{\transnote}[1]{%

  \par\smallskip\noindent\textcolor{notecolor}{\small[Translator's Note: #1]}\par\smallskip

}



%% -------- Hyperref Setup --------

\hypersetup{

  colorlinks=true,

  linkcolor=sectioncolor,

  citecolor=sectioncolor,

  urlcolor=sectioncolor,

  pdfencoding=auto,

  pdftitle={Suanxue Qimeng - Bilingual Edition},

  pdfauthor={Zhu Shijie / Translator},

}



%% ========================================================================

\begin{document}



%% ========================================================================

%% TITLE PAGE

%% ========================================================================

\begin{titlepage}

\begin{center}

\vspace*{1.5cm}



{\fontsize{22}{28}\selectfont\color{titlecolor}

  \textbf{Bilingual Edition}}



\vspace{0.5cm}



{\fontsize{28}{36}\selectfont\color{titlecolor}

  新編算學啓蒙 · Suanxue Qimeng}



\vspace{0.5cm}



{\fontsize{16}{22}\selectfont\color{sectioncolor}

  Introduction to Mathematical Studies}



\vspace{0.3cm}



{\large\color{sectioncolor}

  English--Chinese Parallel Text · 英漢對照}



\vspace{1.5cm}



{\Large\color{titlecolor} 卷上 · Volume I (Part I)}



\vspace{0.3cm}



{\normalsize General Principles (總括) and Volume I Chapters 1--8}



\vspace{2cm}



{\large 元 · 朱世傑撰}



\vspace{0.2cm}



{\normalsize\color{sectioncolor}

  By Zhu Shijie (c.\ 1260--1320 CE)}



\vspace{0.2cm}



{\small\color{sectioncolor}

  Yuan Dynasty, written in 1299 (大德己亥)}



\vspace{2cm}



{\small\color{notecolor}

  Bilingual edition with parallel Chinese-English columns\\

  Preserving the 問/答/術 (Problem/Answer/Method) structure\\

  With translator's notes on rod numeral notation and technical terms

}



\vspace{0.5cm}



{\small\color{notecolor}

  Modern LaTeX edition\ using XeLaTeX and xeCJK\\

  Microsoft YaHei · Times New Roman

}



\end{center}

\end{titlepage}



\tableofcontents

\newpage





%% ========================================================================

%% PREFACE

%% ========================================================================

\bilingualsection{Preface}{序}



\bilingual{%

  算學啓蒙者，元朱世傑之所著也。世傑字漢卿，號松庭，燕山人。

  是書成於大德己亥（1299年），凡三卷，二十門，二百五十九問。

  上卷總括、縱橫因法、身外加法、留頭乘法、身外減法、九歸除法、

  異乘同除、庫務解稅、折變互差；中卷田畝計量、堆積求和、

  約分通分、開方術；下卷方程正負、天元如積。

  其書由淺入深，條理井然，為初學算學之津梁也。

  此書在中國久已失傳，賴朝鮮刻本得以不墜，清羅士琳於揚州重刻之（1839年）。

  其學東傳日本，影響深遠。

}{%

  The \textit{Suanxue Qimeng} (算學啓蒙, ``Introduction to Mathematical Studies'')

  was written by Zhu Shijie (朱世傑), styled Hanqing (漢卿), sobriquet Songting

  (松庭), of Yanshan (燕山), in 1299 CE (the Dade era, 大德己亥) of the Yuan Dynasty.

  It is an elementary mathematical textbook in three volumes (上中下三卷), containing

  20 chapters (二十門) and 259 problems (凡二百五十九問).



  The text progresses systematically from basic arithmetic operations---multiplication,

  division, addition, and subtraction---through advanced topics such as \textit{tian yuan shu}

  (天元術, the method of the celestial element for solving equations), \textit{duo ji}

  (垛積, pile accumulation / summation of series), and \textit{zhao cha} (招差, finite

  differences).



  The book opens with a \textbf{General Principles} (總括) section containing 18

  fundamental rules and rhymes, followed by Volume I (卷上) with eight chapters

  and 113 problems covering multiplication shortcuts, division methods, and

  proportional reasoning. Many problems reflect the economic realities of Yuan

  dynasty China---taxation, commodity prices, land measurement, and military logistics.



  The \textit{Suanxue Qimeng} was lost in China for centuries until a Korean edition

  was rediscovered. It was republished by Luo Shilin (羅士琳) in Yangzhou in 1839.

  The work profoundly influenced the development of mathematics in both Korea and

  Japan.

}



\transnote{This bilingual edition preserves the original Chinese text alongside

an English translation. Technical terms are rendered with their Chinese

characters on first occurrence. The 問/答/術 structure is preserved as

\textbf{Problem}/\textbf{Answer}/\textbf{Method}.}



\longline



%% ========================================================================

%% GENERAL PRINCIPLES

%% ========================================================================

\bilingualsection{General Principles}{總括}

\label{sec:zongguo}



\bilingual{%

  總括者，算學之總要也。凡乘法口訣、除法歌訣、度量衡制、

  大數小數、正負之術、開方之法，皆具於此，為全書之綱領。

  初學者必熟讀之，以為入門之基也。以下列一十八條。

}{%

  The General Principles section presents the foundational knowledge required for

  the entire treatise. It includes multiplication tables, division rhymes, unit

  conversions, notation rules, and algebraic conventions. The following 18 items

  are listed at the beginning of the book. These are the fundamental rules that

  every student must master before proceeding to the problems.

}



%% ========================================================================

\bilingualsubsection{釋九數法}{The Nine Numbers (Multiplication Table)}

%% ========================================================================



\bilingual{%

  乘法口訣，算學之基本也。九九之數，自一一如一起，至九九八十一止，

  凡四十五句。縱橫因之，無所不通。



  一一如一，一二如二，二二如四。

  一三如三，二三如六，三三如九。

  一四如四，二四如八，三四一十二，四四一十六。

  一五如五，二五一十，三五一十五，四五二十，五五二十五。

  一六如六，二六一十二，三六一十八，四六二十四，五六三十，六六三十六。

  一七如七，二七一十四，三七二十一，四七二十八，五七三十五，六七四十二，七七四十九。

  一八如八，二八一十六，三八二十四，四八三十二，五八四十，六八四十八，七八五十六，八八六十四。

  一九如九，二九一十八，三九二十七，四九三十六，五九四十五，六九五十四，七九六十三，八九七十二，九九八十一。

}{%

  The multiplication table is the foundation of all calculation. These ``nine-nines''

  rhymes run from ``one-one is one'' to ``nine-nine is eighty-one,'' totaling 45 formulas.

  By multiplying vertically and horizontally, all products may be obtained.



  This is the complete multiplication table as memorized by every Chinese student.

  Each line begins with the multiplication by 1, then 2, up to the number itself.

  For example, the line for 7 gives: $1 \times 7 = 7$, $2 \times 7 = 14$, $3 \times 7 = 21$,

  $4 \times 7 = 28$, $5 \times 7 = 35$, $6 \times 7 = 42$, $7 \times 7 = 49$.

}



\transnote{The multiplication rhymes follow a pattern: ``$a$ $b$ like $c$'' means $a \times b = c$.

When the product is ten or more, the word ``ten'' (十) is inserted: 三四\textbf{一十}二 means

$3 \times 4 = 12$.}



%% ========================================================================

\bilingualsubsection{九歸除法}{The Nine Returns (Division Rhymes)}

%% ========================================================================



\bilingual{%

  \textit{按古法多用商除。爲初學者難入，則後人以此法代之，卽非正術也。}

  九歸者，以一位數除多位數之歌訣也。其法用歸訣立商，以次減實，

  實盡為度。初學者先習此訣，而後可入商除之門。



  一歸如一進，見一進成十。

  二一添作五，逢二進成十。

  三一三十一，三二六十二，逢三進成十。

  四一二十二，四二添作五，四三七十二，逢四進成十。

  五歸添一倍，逢五進成十。

  六一下加四，六二三十二，六三添作五，六四六十四，六五八十二，逢六進成十。

  七一下加三，七二下加六，七三四十二，七四五十五，七五七十一，七六八十四，逢七進成十。

  八一下加二，八二下加四，八三下加六，八四添作五，八五六十二，八六七十四，八七八十六，逢八進成十。

  九歸隨身下，逢九進成十。

}{%

  \textit{According to ancient methods, the standard division (商除) was commonly used.

  But since beginners find it difficult, later generations substituted this method,

  even though it is not the orthodox technique.}



  The ``nine returns'' are rhymes for dividing a multi-digit number by a single digit.

  The method uses the return-rhymes to establish the quotient digit, then subtracts

  from the dividend repeatedly until the dividend is exhausted.



  \textbf{1-return:} One-return is like one-advance; see one, advance to make ten.\\

  \textbf{2-return:} Two-one add to make five; encounter two, advance to make ten.\\

  \textbf{3-return:} Three-one: thirty-one; three-two: sixty-two; encounter three, advance to make ten.\\

  \textbf{4-return:} Four-one: twenty-two; four-two add to make five; four-three: seventy-two; encounter four, advance to make ten.\\

  \textbf{5-return:} Five-return add one-fold; encounter five, advance to make ten.\\

  \textbf{6-return:} Six-one: add four below; six-two: thirty-two; six-three add to make five; six-four: sixty-four; six-five: eighty-two; encounter six, advance to make ten.\\

  \textbf{7-return:} Seven-one: add three below; seven-two: add six below; seven-three: forty-two; seven-four: fifty-five; seven-five: seventy-one; seven-six: eighty-four; encounter seven, advance to make ten.\\

  \textbf{8-return:} Eight-one: add two below; eight-two: add four below; eight-three: add six below; eight-four add to make five; eight-five: sixty-two; eight-six: seventy-four; eight-seven: eighty-six; encounter eight, advance to make ten.\\

  \textbf{9-return:} Nine-return follows the body down; encounter nine, advance to make ten.

}



\transnote{These rhymes encode single-digit division rules for use on the counting board.

``添作五'' (add to make five) means the quotient digit is 5. ``下加$n$'' means add $n$

to the digit below (i.e., the quotient digit is 1 with remainder $n$). ``逢$N$進成十''

means ``when encountering $N$ or more, place 1 in the next higher position.''}



%% ========================================================================

\bilingualsubsection{斤下留法}{Converting Taels to Jin}

%% ========================================================================



\bilingual{%

  \textit{斤下帶兩者，當以十六約之。今則就省，以此代之也。}

  十六兩為一斤，故兩化斤，須以十六除之。以下所列，即兩數化為斤之小數也。



  一退六二五，二留一二五，三留一八七五，四留二五，五留三一二五，

  六留三七五，七留四三七五，八留單五，九留五六二五，十留六二五，

  十一留六八七五，十二留七五，十三留八一二五，十四留八七五，十五留九三七五。

}{%

  \textit{When jin (斤) has taels (兩) below, one should divide by 16. This is a simplified

  method for that purpose.}



  Since 16 taels = 1 jin, converting taels to jin requires division by 16.

  The following are the decimal fractions for each number of taels:\\

  1 tael = 0.0625 jin, 2 taels = 0.125 jin, 3 taels = 0.1875 jin,

  4 taels = 0.25 jin, 5 taels = 0.3125 jin, 6 taels = 0.375 jin,

  7 taels = 0.4375 jin, 8 taels = 0.5 jin, 9 taels = 0.5625 jin,

  10 taels = 0.625 jin, 11 taels = 0.6875 jin, 12 taels = 0.75 jin,

  13 taels = 0.8125 jin, 14 taels = 0.875 jin, 15 taels = 0.9375 jin.

}





%% ========================================================================

\bilingualsubsection{明縱橫訣}{Rules for Counting Rod Notation}

%% ========================================================================



\bilingual{%

  算籌記數之法，縱橫相間，以別位值。其訣曰：



  一縱十橫，百立千僵。

  千十相望，萬百相當。

  滿六已上，五在上方。

  六不積聚，五不單張。

  言十自過，不滿自當。

  若明此訣，可習九章。

}{%

  The counting rod (籌算) notation system uses alternating vertical and horizontal

  orientations to distinguish place values. The rules state:



  \textbf{Units are vertical, tens are horizontal, hundreds stand upright, thousands lie flat.}\\

  \textbf{Thousands and tens face each other; myriads (10,000s) and hundreds match.}\\

  \textbf{For six and above, the five-bar sits above.}\\

  \textbf{Six does not pile up; five does not stand alone.}\\

  \textbf{When reaching ten, carry over; when less than ten, keep as-is.}\\

  \textbf{If you understand this rhyme, you may study the Nine Chapters.}



  In rod numerals, digits 1--4 are represented by 1--4 vertical (or horizontal) rods.

  The digit 5 is a single horizontal bar (for units) or vertical bar (for tens).

  Digits 6--9 combine the 5-bar with 1--4 additional rods below it.

}



\transnote{The rod numeral system is a decimal place-value notation used on a counting board.

Alternating orientations (vertical/horizontal) between adjacent digits prevents confusion.}



%% ========================================================================

\bilingualsubsection{大數之類}{Large Numbers}

%% ========================================================================



\bilingual{%

  \textit{凡數之大者，天莫能蓋、地莫能載，其數不能極，故謂之大數也。}

  中國大數之名，以萬萬遞進，與西洋之以千遞進者不同。



  一、十、百、千、萬、十萬、百萬、千萬，萬萬曰億，

  萬萬億曰兆，萬萬兆曰京，萬萬京曰陔，萬萬陔曰秭，

  萬萬秭曰壤，萬萬壤曰溝，萬萬溝曰澗，萬萬澗曰正，

  萬萬正曰載，萬萬載曰極，萬萬極曰恆河沙，

  萬萬恆河沙曰阿僧祗，萬萬阿僧祗曰那由他，

  萬萬那由他曰不可思議，萬萬不可思議曰無量數。

}{%

  \textit{As for large numbers: heaven cannot cover them, earth cannot contain them;

  their extent is limitless, hence they are called ``large numbers.''}



  Chinese large number names progress by powers of $10^8$ (萬萬), unlike the Western

  system which progresses by powers of $10^3$ (thousands):



  One, ten, hundred, thousand, myriad ($10^4$), $10^5$, $10^6$, $10^7$,

  $10^8$ = yi (億), $10^{16}$ = zhao (兆), $10^{24}$ = jing (京),

  $10^{32}$ = gai (陔), $10^{40}$ = zi (秭), $10^{48}$ = rang (壤),

  $10^{56}$ = gou (溝), $10^{64}$ = jian (澗), $10^{72}$ = zheng (正),

  $10^{80}$ = zai (載), $10^{88}$ = ji (極), $10^{96}$ = Ganges sand (恆河沙),

  $10^{104}$ = asa.mkhyeya (阿僧祗), $10^{112}$ = nayuta (那由他),

  $10^{120}$ = unthinkable (不可思議), $10^{128}$ = immeasurable (無量數).

}



%% ========================================================================

\bilingualsubsection{小數之類}{Small Numbers}

%% ========================================================================



\bilingual{%

  \textit{凡數之小者，視之無形、取之無像，數亦不能盡，故謂之小數也。}

  小數之名，亦以萬萬遞退。



  一、分、釐、毫、絲、忽、微、纖、沙，萬萬塵曰沙，

  萬萬埃曰塵，萬萬渺曰埃，萬萬漠曰渺，萬萬模糊曰漠，

  萬萬逡巡曰模糊，萬萬須臾曰逡巡，萬萬瞬息曰須臾，

  萬萬彈指曰瞬息，萬萬剎那曰彈指，萬萬六德曰剎那，

  萬萬虛曰六德，萬萬空曰虛，萬萬清曰空，萬萬淨曰清。

  千萬淨，百萬淨，十萬淨，萬淨，千淨，百淨，十淨，一淨。

}{%

  \textit{As for small numbers: looking at them, they have no form; grasping them, they have

  no image; their fineness cannot be exhausted, hence they are called ``small numbers.''}



  Small number names also descend by powers of $10^{-8}$:



  One (1), fen (分, $10^{-1}$), li (釐, $10^{-2}$), hao (毫, $10^{-3}$),

  si (絲, $10^{-4}$), hu (忽, $10^{-5}$), wei (微, $10^{-6}$),

  xian (纖, $10^{-7}$), sha (沙, $10^{-8}$), then continuing:

  $10^{-16}$ = dust (塵), $10^{-24}$ = mote (埃),

  $10^{-32}$ = minute (渺), $10^{-40}$ = desert (漠),

  $10^{-48}$ = blur (模糊), $10^{-56}$ = hesitation (逡巡),

  $10^{-64}$ = instant (須臾), $10^{-72}$ = blink (瞬息),

  $10^{-80}$ = finger-snap (彈指), $10^{-88}$ = k.sana (剎那),

  $10^{-96}$ = liude (六德), $10^{-104}$ = void (虛),

  $10^{-112}$ = emptiness (空), $10^{-120}$ = clarity (清),

  $10^{-128}$ = purity (淨).

}



%% ========================================================================

\bilingualsubsection{求諸率類}{Conversion Rates}

%% ========================================================================



\bilingual{%

  諸物互求之法，以乘除為用。其率如下：



  兩求銖二十四乘，銖求兩二十四除。

  斤求兩身外加六，兩求斤身外減六。

  秤求斤身外加五，斤求秤身外減五。

  據物賣錢而用乘，據錢買物而用除。

}{%

  To convert between units, use multiplication and division with these rates:



  Taels (兩) to zhu (銖): multiply by 24. Zhu to taels: divide by 24.

  (1 tael = 24 zhu)



  Jin (斤) to taels: multiply by 16 (``add six to the body,''

  i.e., the conversion factor is 16 = 10 + 6). Taels to jin: divide by 16.



  Cheng (秤) to jin: multiply by 15 (``add five to the body,''

  since 1 cheng = 15 jin). Jin to cheng: divide by 15.



  To sell goods for money: use multiplication.

  To buy goods with money: use division.

}



%% ========================================================================

\bilingualsubsection{斛㪷起率}{Volume Measurements}

%% ========================================================================



\bilingual{%

  量之起率，始於微細，積而為大。



  量起於圭（六粒之粟）。十圭謂之一撮，十撮謂之一抄，

  十抄謂之一勺，十勺謂之一合，十合謂之一升，

  十升謂之一㪷，十㪷謂之一斛。

}{%

  Volume measurement begins from the smallest unit and accumulates upward:



  Volume begins with the gui (圭): six grains of millet.

  10 gui = 1 cuo (撮), 10 cuo = 1 chao (抄),

  10 chao = 1 shao (勺), 10 shao = 1 ge (合),

  10 ge = 1 sheng (升), 10 sheng = 1 dou (㪷),

  10 dou = 1 hu (斛).



  Thus 1 hu = $10^7$ gui.

}



%% ========================================================================

\bilingualsubsection{斤秤起率}{Weight Measurements}

%% ========================================================================



\bilingual{%

  衡之起率，以黍為基。



  衡起於黍（形大如粟）。十黍謂之一絫，十絫謂之一銖，

  六銖謂之一分，四分謂之一兩，十六兩謂之一斤，

  十五斤謂之一秤，三十斤謂之一鈞，四鈞謂之一碩（重一百二十斤）。

}{%

  Weight measurement begins from the smallest grain:



  Weight begins with the shu (黍): a grain the size of millet.

  10 shu = 1 lei (絫), 10 lei = 1 zhu (銖),

  6 zhu = 1 fen (分), 4 fen = 1 liang (兩, tael),

  16 liang = 1 jin (斤), 15 jin = 1 cheng (秤),

  30 jin = 1 jun (鈞), 4 jun = 1 shuo (碩, = 120 jin).

}



%% ========================================================================

\bilingualsubsection{端匹起率}{Length Measurements}

%% ========================================================================



\bilingual{%

  度之起率，以蠶絲為始。



  度起於忽（蠶吐之絲）。十忽謂之一絲，十絲謂之一毫，

  十毫謂之一釐，十釐謂之一分，十分謂之一寸，

  十寸謂之一尺，十尺謂之一丈。

  匹率（或三丈二，或二丈四）。

  端率（或五丈五，或四丈八）。

}{%

  Length measurement begins from the finest silk thread:



  Length begins with the hu (忽): a single silk thread from a silkworm.

  10 hu = 1 si (絲), 10 si = 1 hao (毫),

  10 hao = 1 li (釐), 10 li = 1 fen (分),

  10 fen = 1 cun (寸, inch), 10 cun = 1 chi (尺, foot),

  10 chi = 1 zhang (丈).



  Cloth bolt (匹) measure: either 3.2 zhang or 2.4 zhang.\\

  Cloth roll (端) measure: either 5.5 zhang or 4.8 zhang.

}



%% ========================================================================

\bilingualsubsection{田畝起率}{Land Area Measurements}

%% ========================================================================



\bilingual{%

  田畝之制，以步為基。



  田起於忽（闊一寸，長六寸）。十忽謂之一絲，十絲謂之一毫，

  十毫謂之一釐，十釐謂之一分，十分謂之一畝，百畝謂之一頃。

  三百步謂之一里。



  \textit{按畝法，闊一步，長二百四十步，當自方五尺為步也。

  其里法，三百步為里者，當自方六尺為步。若三百六十步為里者，

  當以自方五尺為步也。}

}{%

  The land area system uses the ``pace'' (步) as its foundation:



  Land area begins with the hu (忽): 1 cun wide by 6 cun long.

  10 hu = 1 si (絲), 10 si = 1 hao (毫),

  10 hao = 1 li (釐), 10 li = 1 fen (分),

  10 fen = 1 mu (畝), 100 mu = 1 qing (頃).

  300 paces = 1 li (里).



  \textit{According to the mu method: 1 pace wide by 240 paces long = 1 mu,

  using a 5-chi pace. For the li method: 300 paces = 1 li using a 6-chi pace;

  or 360 paces = 1 li using a 5-chi pace.}



  One mu $\approx$ 1/15 of a hectare $\approx$ 667 m$^2$ in modern terms.

}





%% ========================================================================

\bilingualsubsection{古法圓率}{Ancient Circle Ratio}

%% ========================================================================



\bilingual{%

  古之圓率，以周三徑一為率。



  周三尺，徑一尺。

}{%

  The ancient ratio for the circle:



  Circumference 3 chi, diameter 1 chi.



  This gives $\pi = 3$, the ancient ratio ``three for circumference, one for diameter.''

  This value was used in the \textit{Zhoubi Suanjing} (周髀算經) and the

  \textit{Jiuzhang Suanshu} (九章算術) before more precise values were obtained.

}



%% ========================================================================

\bilingualsubsection{劉徽新術}{Liu Hui's New Method}

%% ========================================================================



\bilingual{%

  \textit{劉徽乃魏人也，立此新術以究圓之幽微。}

  劉徽以割圓術求圓周率，割圓為多邊形，邊數愈多則愈近於圓。



  周一百五十七尺，徑五十尺。

}{%

  \textit{Liu Hui was a man of Wei; he established this new method to investigate

  the subtle mysteries of the circle.}



  Liu Hui (fl.\ 263 CE) used the ``method of cutting the circle'' (割圓術):

  inscribing polygons with increasingly many sides to approximate the circle.

  He calculated up to a 3072-sided polygon.



  Circumference 157 chi, diameter 50 chi.



  This gives $\pi = \dfrac{157}{50} = 3.14$, a significant improvement over $\pi = 3$.

}



%% ========================================================================

\bilingualsubsection{沖之密率}{Zu Chongzhi's Dense Ratio}

%% ========================================================================



\bilingual{%

  \textit{沖之姓祖，乃宋南徐州從事史。立此密率亦究圓之微也。}

  祖沖之更以密率求之，其精密為當時之最。



  周二十二尺，徑七尺。

}{%

  \textit{Chongzhi, surnamed Zu, was an officer of Nanku province in the Song dynasty.

  He established this dense ratio also to investigate the circle's subtlety.}



  Zu Chongzhi (429--500 CE) gave two approximations for $\pi$:

  the ``approximate ratio'' (約率) $\pi \approx 22/7$ and

  the ``dense ratio'' (密率) $\pi \approx 355/113$.

  The dense ratio gives $\pi \approx 3.1415929\ldots$, accurate to six decimal places.



  Circumference 22 chi, diameter 7 chi.



  This gives $\pi = \dfrac{22}{7} \approx 3.142857$ (the ``approximate ratio'').

}



\transnote{Zu Chongzhi's dense ratio $\frac{355}{113}$ is one of the best rational

approximations of $\pi$ with small denominator. It was not rediscovered in Europe

until Adriaan Anthoniszoon (c.\ 1600).}



%% ========================================================================

\bilingualsubsection{明異名訣}{Names of Fractions}

%% ========================================================================



\bilingual{%

  \textit{此乃劇漏之數以巧呼之。}

  古之算家，以巧名分數，以便口訣之用。



  二分之一為中半，三分之一為少半，三分之二為太半，

  四分之一為弱半，四分之三為強半。

}{%

  \textit{These are special names for fractions, called by elegant names.}



  Ancient mathematicians gave special names to common fractions for ease of use:



  $\dfrac{1}{2}$ = ``middle half'' (中半),\quad

  $\dfrac{1}{3}$ = ``lesser half'' (少半),\quad

  $\dfrac{2}{3}$ = ``greater half'' (太半),\\

  $\dfrac{1}{4}$ = ``weak half'' (弱半),\quad

  $\dfrac{3}{4}$ = ``strong half'' (強半).

}



%% ========================================================================

\bilingualsubsection{明正負術}{Rules for Positive and Negative Numbers}

%% ========================================================================



\bilingual{%

  正負之術，為方程之根本。同名相減，異名相加，

  正無入負之，負無入正之。



  其同名相減，則異名相加；正無入負之，負無入正之。

  其異名相減，則同名相加；正無入正之，負無入負之。



  \textit{按九章注云：兩算得失相返，要令正負以名之。

  正算赤、負算黑，不則以邪正為異。其無入者，為無對也。

  無所得減，則使消奪者居位也。（「人」作「入」非。）}

}{%

  The rules for positive and negative numbers are fundamental to solving systems

  of equations. The rules state:



  \textbf{Same signs subtract, different signs add.} When subtracting a positive

  from zero (or a smaller positive), the result is negative; when subtracting a

  negative from zero, the result is positive.



  \textbf{Rule 1:} Same signs subtract $\rightarrow$ different signs add;

  positive with nothing enters negative; negative with nothing enters positive.



  \textbf{Rule 2:} Different signs subtract $\rightarrow$ same signs add;

  positive with nothing enters positive; negative with nothing enters negative.



  \textit{According to the commentary on the Jiuzhang: two calculations with opposite

  outcomes are named positive and negative. Positive rods are red, negative rods are black;

  otherwise, slanted versus straight distinguishes them. ``Nothing to enter'' means there

  is no counterpart; when there is nothing to subtract from, the consuming quantity

  takes the position.}



  These rules correspond to: $(+a) - (+b) = +(a-b)$, $(-a) - (-b) = -(a-b)$,

  $(+a) + (-b) = +(a-b)$, $(+a) + (-b) = -(b-a)$ when $b > a$, etc.

}



\transnote{These are among the earliest known rules for signed arithmetic in world mathematics.

The rod-color convention (red = positive, black = negative) dates to the Han dynasty.

These rules predate European acceptance of negative numbers by over a millennium.}



%% ========================================================================

\bilingualsubsection{明乘除段}{Rules for Multiplication and Division}

%% ========================================================================



\bilingual{%

  乘除之術，凡幾何名義，不可不辨。



  長平相併曰和，長平相減曰較。

  長平相乘曰積，自相稱之曰冪。

  同名相乘曰正，異名相乘為負。

  平除長為小長，長除平為小平。

  小長平相併曰小和，小長平相減餘小較，小長平相乘得一步為小積。

}{%

  The terminology of multiplication and division must be clearly understood:



  Length plus width is called ``harmony'' (和, sum).\\

  Length minus width is called ``comparison'' (較, difference).\\

  Length times width is called ``accumulation'' (積, area/product).\\

  A number multiplied by itself is called ``power'' (冪, square).\\

  Same signs multiplied give positive; different signs give negative.\\

  Width divided by length gives ``small length''; length divided by width gives ``small width.''\\

  Small length plus small width = ``small harmony''; their difference = ``small comparison'';

  their product = ``small accumulation.''

}



\transnote{The terms 長 (long) and 平 (broad/width) denote the sides of a rectangle.

These algebraic conventions were essential for the tian yuan shu (天元術) method

that Zhu Shijie would develop in his later work, the \textit{Siyuan Yujian}.}



%% ========================================================================

\bilingualsubsection{明開方法}{Rules for Root Extraction}

%% ========================================================================



\bilingual{%

  開方之術，方程之基也。凡開平方、開立方，皆用此術。



  置積為實，及方廉隅，同加異減開之。

}{%

  The method of root extraction is the foundation of solving equations.

  It applies to square roots, cube roots, and higher roots.



  Set up the accumulation (積) as the dividend (實), together with the side (方),

  edge (廉), and corner (隅). Like signs add, unlike signs subtract, and extract.



  In the root extraction process for $x^2 = N$, the terms are:

  \textbf{shi} (實, dividend) = $N$,

  \textbf{fang} (方, side) = the linear term coefficient,

  \textbf{lian} (廉, edge) = the edge term,

  \textbf{yu} (隅, corner) = the leading coefficient.

  The ``incremental root extraction'' (增乘開方) method systematically

  refines the root digit by digit.

}



\transnote{The ``incremental multiplication'' method (增乘開方), attributed to Jia Xian

(11th century), is a precursor to what is now known as the Ruffini-Horner method

for polynomial root-finding. It was systematized by Zhu Shijie in his later work.}



\longline





%% ========================================================================

%% VOLUME I

%% ========================================================================

\bilingualsection{Volume I}{卷上}

\label{sec:volume1}



\bilingual{%

  卷上凡八門，一百一十三問。其門曰：縱橫因法、身外加法、

  留頭乘法、身外減法、九歸除法、異乘同除、庫務解稅、折變互差。

  大要皆乘除之術，為算學之初步也。

}{%

  Volume I contains eight chapters (八門) with 113 problems (一百一十三問),

  covering multiplication shortcuts, division methods, and proportional reasoning.

  Many problems reflect the economic realities of Yuan dynasty China---taxation,

  commodity prices, land measurement, and military logistics.

}



\longline



%% ========================================================================

%% Chapter 1: 縱橫因法門

%% ========================================================================

\bilingualsubsection{縱橫因法門}{The Vertical--Horizontal Multiplication Method}

\label{sec:zongheng}



\bilingual{%

  \begin{rhymecn}

  此法從來向上因，但言十者過其身，\\

  呼如本位須當作，知算縱橫數目真。

  \end{rhymecn}

  縱橫因法者，乘法之初門也。其法列實於上，以法數從上因之，

  言十過身，言如下位，逐位相乘，得數即積也。凡八問。

}{%

  \begin{rhymecn}

  This method always multiplies upward from below,\\

  When ten is reached, carry past the body;\\

  Call ``as-is'' for the current place value;\\

  Know the vertical--horizontal counts to be true.

  \end{rhymecn}

  This rhyme describes the ``vertical--horizontal'' multiplication method:

  multiply from bottom to top; when the product is ten or more, carry to the

  next position; call ``as-is'' for the current place value. This is the most

  basic multiplication technique.

}



\bilingual{%

  \textbf{問1.}

  今有粟二百一十六斗，每斗價錢二文。問計錢幾何？\\

  \textbf{答曰：}四百三十二文。\\

  \textbf{術曰：}列物數在上，各以價錢從上因之，即得。

}{%

  \textbf{Problem 1.}

  Suppose there are 216 dou of millet, each dou costing 2 cash.

  How much cash in total?\\

  \textbf{Answer:} 432 cash.\\

  \textbf{Method:} Place the quantity above; multiply each by the unit price

  from above downward; the result is obtained.\\

  (Calculation: $216 \times 2 = 432$)

}



\bilingual{%

  \textbf{問2.}

  今有絲一百四十四兩，每兩價錢三百文。問計錢幾何？\\

  \textbf{答曰：}四十三貫二百文。\\

  \textbf{術曰：}列物數在上，各以價錢從上因之，即得。

}{%

  \textbf{Problem 2.}

  Suppose there are 144 liang of silk, each liang costing 300 cash.

  How much cash in total?\\

  \textbf{Answer:} 43 guan 200 cash.\\

  \textbf{Method:} Place the quantity above; multiply each by the unit price; obtain.\\

  (Calculation: $144 \times 300 = 43{,}200$)

}



\bilingual{%

  \textbf{問3.}

  今有羊三百五十四只，每只價錢四貫文。問計錢幾何？\\

  \textbf{答曰：}一千四百一十六貫。\\

  \textbf{術曰：}列物數在上，各以價錢從上因之，即得。

}{%

  \textbf{Problem 3.}

  Suppose there are 354 sheep, each costing 4 guan.

  How much cash in total?\\

  \textbf{Answer:} 1,416 guan.\\

  \textbf{Method:} Place the quantity above; multiply by the unit price; obtain.\\

  (Calculation: $354 \times 4 = 1{,}416$)

}



\bilingual{%

  \textbf{問4.}

  今有銀五百四十七錠，每錠重五十兩。問為兩幾何？\\

  \textbf{答曰：}二萬七千三百五十兩。\\

  \textbf{術曰：}列物數在上，各以錠重從上因之，即得。

}{%

  \textbf{Problem 4.}

  Suppose there are 547 ingots of silver, each ingot weighing 50 liang.

  How many liang in total?\\

  \textbf{Answer:} 27,350 liang.\\

  \textbf{Method:} Place the quantity above; multiply each by the ingot weight; obtain.\\

  (Calculation: $547 \times 50 = 27{,}350$)

}



\bilingual{%

  \textbf{問5.}

  今有絹七百三十六匹，每匹價錢六貫文。問計錢幾何？\\

  \textbf{答曰：}四千四百一十六貫。\\

  \textbf{術曰：}列物數在上，各以價錢從上因之，即得。

}{%

  \textbf{Problem 5.}

  Suppose there are 736 bolts of silk, each costing 6 guan.

  How much cash in total?\\

  \textbf{Answer:} 4,416 guan.\\

  \textbf{Method:} Place the quantity above; multiply by unit price; obtain.\\

  (Calculation: $736 \times 6 = 4{,}416$)

}



\bilingual{%

  \textbf{問6.}

  今有麻八百九十二秤，每秤價錢七百文。問計錢幾何？\\

  \textbf{答曰：}六百二十四貫四百文。\\

  \textbf{術曰：}列物數在上，各以價錢從上因之，即得。

}{%

  \textbf{Problem 6.}

  Suppose there are 892 cheng of hemp, each cheng costing 700 cash.

  How much cash in total?\\

  \textbf{Answer:} 624 guan 400 cash.\\

  \textbf{Method:} Place the quantity above; multiply by unit price; obtain.\\

  (Calculation: $892 \times 700 = 624{,}400$)

}



\bilingual{%

  \textbf{問7.}

  今有布六百三十四尺，每尺價錢八十文。問計錢幾何？\\

  \textbf{答曰：}五十貫七百二十文。\\

  \textbf{術曰：}列物數在上，各以價錢從上因之，即得。

}{%

  \textbf{Problem 7.}

  Suppose there are 634 chi of cloth, each chi costing 80 cash.

  How much cash in total?\\

  \textbf{Answer:} 50 guan 720 cash.\\

  \textbf{Method:} Place the quantity above; multiply by unit price; obtain.\\

  (Calculation: $634 \times 80 = 50{,}720$)

}



\bilingual{%

  \textbf{問8.}

  今有馬四百二十五匹，每匹價錢九十貫。問計錢幾何？\\

  \textbf{答曰：}三萬八千二百五十貫。\\

  \textbf{術曰：}列物數在上，各以價錢從上因之，即得。

}{%

  \textbf{Problem 8.}

  Suppose there are 425 horses, each costing 90 guan.

  How much cash in total?\\

  \textbf{Answer:} 38,250 guan.\\

  \textbf{Method:} Place the quantity above; multiply by unit price; obtain.\\

  (Calculation: $425 \times 90 = 38{,}250$)

}



\longline



%% ========================================================================

%% Chapter 2: 身外加法門

%% ========================================================================

\bilingualsubsection{身外加法門}{Addition to the Body}

\label{sec:shenwaijia}



\bilingual{%

  \begin{rhymecn}

  算中加法最堪誇，言十之時就位加，\\

  但遇呼如身下列，君從法式定無差。

  \end{rhymecn}

  身外加法者，乘法之捷法也。其法以實為身，於身外逐位加之，

  言十就身，言如下列。凡十問。

}{%

  \begin{rhymecn}

  Among calculations, the addition method is most praiseworthy;\\

  When ten is reached, add at the current place;\\

  When ``as-is'' is called, place it below the body;\\

  Follow this rule and you will make no error.

  \end{rhymecn}

  This rhyme describes a shortcut multiplication method: add to the

  ``body'' (the multiplicand itself). When the product reaches ten, add at the

  current position; when it is less than ten (``as-is''), place it below the

  current digit. This is essentially a streamlined single-digit multiplication

  performed on the counting board.

}



\bilingual{%

  \textbf{問9.}

  今有米六碩八斗四升，每斗價錢一百一十文。問計錢幾何？\\

  \textbf{答曰：}七貫五百二十四文。\\

  \textbf{術曰：}列物數在上，以價錢從上身外加之，即得。

}{%

  \textbf{Problem 9.}

  Suppose there are 6 shuo 8 dou 4 sheng of rice, each dou costing 110 cash.

  How much cash in total?\\

  \textbf{Answer:} 7 guan 524 cash.\\

  \textbf{Method:} Place the quantity above; add the price to the body from above; obtain.

}



\bilingual{%

  \textbf{問10.}

  今有羅三十四尺六寸，每尺價錢一百二十文。問計錢幾何？\\

  \textbf{答曰：}四貫一百五十二文。\\

  \textbf{術曰：}列物數在上，以價錢從上身外加之，即得。

}{%

  \textbf{Problem 10.}

  Suppose there are 34 chi 6 cun of gauze, each chi costing 120 cash.

  How much cash in total?\\

  \textbf{Answer:} 4 guan 152 cash.\\

  \textbf{Method:} Place the quantity above; add to the body; obtain.

}



\bilingual{%

  \textbf{問11.}

  今有鹽八百七十三袋，每袋價錢一十三貫文。問計錢幾何？\\

  \textbf{答曰：}一萬一千三百四十九貫。\\

  \textbf{術曰：}列物數在上，以價錢從上身外加之，即得。

}{%

  \textbf{Problem 11.}

  Suppose there are 873 bags of salt, each bag costing 13 guan.

  How much cash in total?\\

  \textbf{Answer:} 11,349 guan.\\

  \textbf{Method:} Place the quantity above; add to the body; obtain.\\

  (Calculation: $873 \times 13 = 11{,}349$)

}



\bilingual{%

  \textbf{問12.}

  今有麥二十三碩四斗，每斗價錢一百三十文。問計錢幾何？\\

  \textbf{答曰：}三十貫四百二十文。\\

  \textbf{術曰：}列麥數在上，以價錢從上身外加之，即得。

}{%

  \textbf{Problem 12.}

  Suppose there are 23 shuo 4 dou of wheat, each dou costing 130 cash.

  How much cash in total?\\

  \textbf{Answer:} 30 guan 420 cash.\\

  \textbf{Method:} Place the wheat quantity above; add to the body; obtain.

}



\bilingual{%

  \textbf{問13.}

  今有絲四十二斤三兩，每兩價錢二百四十文。問計錢幾何？\\

  \textbf{答曰：}一十貫一百七十五文。\\

  \textbf{術曰：}列絲數在上，以價錢從上身外加之，即得。

}{%

  \textbf{Problem 13.}

  Suppose there are 42 jin 3 liang of silk, each liang costing 240 cash.

  How much cash in total?\\

  \textbf{Answer:} 10 guan 175 cash.\\

  \textbf{Method:} Place the silk quantity above; add to the body; obtain.

}



\bilingual{%

  \textbf{問14.}

  今有粟五十六碩七斗，每斗價錢一百五十文。問計錢幾何？\\

  \textbf{答曰：}八十五貫五十文。\\

  \textbf{術曰：}列粟數在上，以價錢從上身外加之，即得。

}{%

  \textbf{Problem 14.}

  Suppose there are 56 shuo 7 dou of millet, each dou costing 150 cash.

  How much cash in total?\\

  \textbf{Answer:} 85 guan 50 cash.\\

  \textbf{Method:} Place the millet quantity above; add to the body; obtain.\\

  (Calculation: $567 \times 150 = 85{,}050$)

}



\bilingual{%

  \textbf{問15.}

  今有茶三百六十二袋，每袋價錢二貫五百文。問計錢幾何？\\

  \textbf{答曰：}九百五貫。\\

  \textbf{術曰：}列茶數在上，以價錢從上身外加之，即得。

}{%

  \textbf{Problem 15.}

  Suppose there are 362 bags of tea, each bag costing 2 guan 500 cash.

  How much cash in total?\\

  \textbf{Answer:} 905 guan.\\

  \textbf{Method:} Place the tea quantity above; add to the body; obtain.\\

  (Calculation: $362 \times 2.5 = 905$)

}



\bilingual{%

  \textbf{問16.}

  今有綿一千二百四十五兩，每兩價錢六十文。問計錢幾何？\\

  \textbf{答曰：}七十四貫七百文。\\

  \textbf{術曰：}列綿數在上，以價錢從上身外加之，即得。

}{%

  \textbf{Problem 16.}

  Suppose there are 1,245 liang of cotton, each liang costing 60 cash.

  How much cash in total?\\

  \textbf{Answer:} 74 guan 700 cash.\\

  \textbf{Method:} Place the cotton quantity above; add to the body; obtain.\\

  (Calculation: $1{,}245 \times 60 = 74{,}700$)

}



\bilingual{%

  \textbf{問17.}

  今有鐵二千八百七十四斤，每斤價錢四十五文。問計錢幾何？\\

  \textbf{答曰：}一百二十九貫三百三十文。\\

  \textbf{術曰：}列鐵數在上，以價錢從上身外加之，即得。

}{%

  \textbf{Problem 17.}

  Suppose there are 2,874 jin of iron, each jin costing 45 cash.

  How much cash in total?\\

  \textbf{Answer:} 129 guan 330 cash.\\

  \textbf{Method:} Place the iron quantity above; add to the body; obtain.\\

  (Calculation: $2{,}874 \times 45 = 129{,}330$)

}



\bilingual{%

  \textbf{問18.}

  今有鉛一千五百六十三斤四兩，每斤價錢八十文。問計錢幾何？\\

  \textbf{答曰：}一百二十五貫六十文。\\

  \textbf{術曰：}列鉛數在上，以價錢從上身外加之，即得。

}{%

  \textbf{Problem 18.}

  Suppose there are 1,563 jin 4 liang of lead, each jin costing 80 cash.

  How much cash in total?\\

  \textbf{Answer:} 125 guan 60 cash.\\

  \textbf{Method:} Place the lead quantity above; add to the body; obtain.

}



\bilingual{%

  \textbf{問19.}

  今有錫八百七十六斤半，每斤價錢一百七十文。問計錢幾何？\\

  \textbf{答曰：}一百四十九貫五百五文。\\

  \textbf{術曰：}列錫數在上，以價錢從上身外加之，即得。

}{%

  \textbf{Problem 19.}

  Suppose there are 876.5 jin of tin, each jin costing 170 cash.

  How much cash in total?\\

  \textbf{Answer:} 149 guan 505 cash.\\

  \textbf{Method:} Place the tin quantity above; add to the body; obtain.

}



\longline





%% ========================================================================

%% Chapter 3: 留頭乘法門

%% ========================================================================

\bilingualsubsection{留頭乘法門}{Multiplication Keeping the First Digit}

\label{sec:liutou}



\bilingual{%

  \begin{rhymecn}

  留頭乘法別規模，起首先從次位呼，\\

  言十靠身如隔位，遍臨頭位破身鋪。

  \end{rhymecn}

  留頭乘法者，乘法之巧術也。其法以法數之首位暫留不動，

  先從次位乘起，言十靠身，如隔位布之，遍臨頭位則破身鋪之。

  凡二十問。

}{%

  \begin{rhymecn}

  The keep-the-head method has a special pattern;\\

  Begin from the second digit, not the first;\\

  When ten, place next to the body; when ``as-is,'' skip a place;\\

  Finally use the head digit, breaking the body.

  \end{rhymecn}

  The ``keep-the-head'' multiplication method: a special technique where

  the first digit of the multiplier is temporarily kept (left in place), and

  multiplication begins from the second digit. When the product reaches ten,

  it is placed next to the body; ``as-is'' results are placed with a skip.

  Finally the head digit itself is used to multiply, ``breaking'' the original

  body. This avoids shifting the multiplicand on the counting board.

}



\bilingual{%

  \textbf{問20.}

  今有白豆八十四斛，每斛價錢二百一十文。問計錢幾何？\\

  \textbf{答曰：}一十七貫六百四十文。\\

  \textbf{術曰：}列豆數於上，以斛價從上次位呼之，言十靠身，如隔位布之，

  遍臨頭位破身，即得。

}{%

  \textbf{Problem 20.}

  Suppose there are 84 hu of white beans, each hu costing 210 cash.

  How much cash in total?\\

  \textbf{Answer:} 17 guan 640 cash.\\

  \textbf{Method:} Place the bean quantity above; call the hu-price from the

  second digit upward; when ten, place next to the body; when ``as-is,'' skip a place;

  when reaching the head digit, break the body; obtain.\\

  (Calculation: $84 \times 210 = 17{,}640$)

}



\bilingual{%

  \textbf{問21.}

  今有胡椒六十三斤四兩，每斤價錢三百八十文。問計錢幾何？\\

  \textbf{答曰：}二十四貫三十五文。\\

  \textbf{術曰：}列椒數於上，以斤價從上次位呼之，言十靠身，如隔位布之，

  遍臨頭位破身，即得。

}{%

  \textbf{Problem 21.}

  Suppose there are 63 jin 4 liang of pepper, each jin costing 380 cash.

  How much cash in total?\\

  \textbf{Answer:} 24 guan 35 cash.\\

  \textbf{Method:} Place the pepper quantity above; call the price from the

  second digit; when ten, place next to body; when ``as-is,'' skip; at the head

  digit, break the body; obtain.

}



\bilingual{%

  \textbf{問22.}

  今有白檀共數斤，下留兩於上，以四貫九十六文乘之。問得幾何？\\

  \textbf{答曰：}二萬貫。\\

  \textbf{術曰：}列白檀共數斤下留兩於上，以四貫九十六文從上次位呼之，

  言十靠身，如隔位布之，遍臨頭位破身，即得。

}{%

  \textbf{Problem 22.}

  Suppose there is a quantity of white sandalwood in jin,

  leaving taels below, multiply by 4 guan 96 cash.

  What is obtained?\\

  \textbf{Answer:} 20,000 guan.\\

  \textbf{Method:} Place the sandalwood quantity with taels below;

  call 4 guan 96 cash from the second digit; when ten, place next to body;

  when ``as-is,'' skip; at head digit, break body; obtain.

}



\bilingual{%

  \textbf{問23.}

  今有黍三千六百六十二碩一斛九合三勺七抄五撮。問為麥幾何？\\

  \textbf{答曰：}三萬斛。\\

  \textbf{術曰：}列黍數於上，以一斛麥八升一合九勺二抄乘之，即得。

}{%

  \textbf{Problem 23.}

  Suppose there are 3,662 shuo 1 hu 9 ge 3 shao 7 chao 5 cuo of millet.

  How much wheat is it equivalent to?\\

  \textbf{Answer:} 30,000 hu.\\

  \textbf{Method:} Place the millet quantity above; multiply by 1 hu wheat

  = 8 sheng 1 ge 9 shao 2 chao; obtain.

}



\bilingual{%

  \textbf{問24.}

  今有粟七萬八千一百二十五碩，每斛為御米八升一合九勺二抄。問共得幾何？\\

  \textbf{答曰：}三千二百七十斛。\\

  \textbf{術曰：}列粟數於上，以御米率從上次位呼之，即得。

}{%

  \textbf{Problem 24.}

  Suppose there are 78,125 shuo of grain, each hu yielding 8 sheng 1 ge

  9 shao 2 chao of imperial rice. How much rice is obtained?\\

  \textbf{Answer:} 3,270 hu.\\

  \textbf{Method:} Place the grain quantity above; call the imperial rice rate

  from the second digit; obtain.

}



\bilingual{%

  \textbf{問25.}

  今有絲六萬一千六百九十八秤。問為斤幾何？\\

  \textbf{答曰：}一百六十九萬一千六百九十八斤。\\

  \textbf{術曰：}列絲數於上，以一秤十六兩為法，先以六因之，再以四因之，

  即得斤數。

}{%

  \textbf{Problem 25.}

  Suppose there are 61,698 cheng of silk. How many jin?\\

  \textbf{Answer:} 1,691,698 jin.\\

  \textbf{Method:} Place the silk quantity above; using 1 cheng = 16 liang

  (= 1 jin), first multiply by 6, then by 4; obtain the jin number.

}



\bilingual{%

  \textbf{問26.}

  今有絲六萬一千六百九十八秤，每斤重一十六兩。問為兩幾何？\\

  \textbf{答曰：}一千六百九十八萬一千六百九十八兩。\\

  \textbf{術曰：}列絲數於上，以每斤十六兩為法，先從次位呼之，遍臨頭位破身，

  即得。

}{%

  \textbf{Problem 26.}

  Suppose there are 61,698 cheng of silk, each jin weighing 16 liang.

  How many liang?\\

  \textbf{Answer:} 16,981,698 liang.\\

  \textbf{Method:} Place the silk quantity above; using 1 jin = 16 liang,

  call from the second digit; when reaching the head digit, break the body; obtain.

}



\bilingual{%

  \textbf{問27.}

  今有錢七貫四百一十二文，欲令一十七人分之。問人得幾何？\\

  \textbf{答曰：}四百三十六文。\\

  \textbf{術曰：}列錢數於上，以一十七人為法，從上次位呼之，即得。

}{%

  \textbf{Problem 27.}

  Suppose there are 7 guan 412 cash, to be divided among 17 people.

  How much per person?\\

  \textbf{Answer:} 436 cash.\\

  \textbf{Method:} Place the cash quantity above; using 17 people as divisor,

  call from the second digit; obtain.\\

  (Calculation: $7{,}412 \div 17 = 436$)

}



\bilingual{%

  \textbf{問28.}

  今有糧五千八百八十六碩，欲給貧難戶，每戶一碩八斗。問給幾戶？\\

  \textbf{答曰：}三千二百七十戶。\\

  \textbf{術曰：}列糧數於上，以每戶一碩八斗為法，從上次位呼之，即得。

}{%

  \textbf{Problem 28.}

  Suppose there are 5,886 shuo of grain, to be distributed to poor households,

  each household receiving 1 shuo 8 dou. How many households?\\

  \textbf{Answer:} 3,270 households.\\

  \textbf{Method:} Place the grain quantity above; using 1 shuo 8 dou per household

  as divisor, call from the second digit; obtain.\\

  (Calculation: $5{,}886 \div 1.8 = 3{,}270$)

}



\bilingual{%

  \textbf{問29.}

  今有錢六十五貫五百五十文，欲買苧絲，每尺價錢一百九十文。問得幾何？\\

  \textbf{答曰：}三百四十五尺。\\

  \textbf{術曰：}列錢數於上，以尺價為法，從上次位呼之，即得。

}{%

  \textbf{Problem 29.}

  Suppose there are 65 guan 550 cash, to buy ramie silk at 190 cash per chi.

  How many chi are obtained?\\

  \textbf{Answer:} 345 chi.\\

  \textbf{Method:} Place the cash quantity above; using the chi-price as divisor,

  call from the second digit; obtain.\\

  (Calculation: $65{,}550 \div 190 = 345$)

}



\bilingual{%

  \textbf{問30.}

  今有錢一百四貫三百七十二文，欲買麻布，每匹價錢一百九十四文。問買布幾何？\\

  \textbf{答曰：}五百三十八匹。\\

  \textbf{術曰：}列錢數於上，以匹價為法，從上次位呼之，即得。

}{%

  \textbf{Problem 30.}

  Suppose there are 104 guan 372 cash, to buy hemp cloth at 194 cash per bolt.

  How many bolts?\\

  \textbf{Answer:} 538 bolts.\\

  \textbf{Method:} Place the cash quantity above; using the bolt-price as divisor,

  call from the second digit; obtain.\\

  (Calculation: $104{,}372 \div 194 = 538$)

}



\bilingual{%

  \textbf{問31.}

  今有錢五千五百五十八貫三百文，欲買松木，每株價錢一貫七百五文。問買木幾何？\\

  \textbf{答曰：}三千二百六十株。\\

  \textbf{術曰：}列錢數於上，以株價為法，從上次位呼之，即得。

}{%

  \textbf{Problem 31.}

  Suppose there are 5,558 guan 300 cash, to buy pine trees at 1 guan 705 cash each.

  How many trees?\\

  \textbf{Answer:} 3,260 trees.\\

  \textbf{Method:} Place the cash quantity above; using the tree-price as divisor,

  call from the second digit; obtain.\\

  (Calculation: $5{,}558{,}300 \div 1{,}705 = 3{,}260$)

}



\bilingual{%

  \textbf{問32.}

  今有芝麻共數於上，以三斤七分五釐乘之。問得幾何？\\

  \textbf{術曰：}列芝麻共數於上，以三斤七分五釐為法，從上次位呼之，

  言十靠身，遍臨頭位破身，即得。

}{%

  \textbf{Problem 32.}

  Suppose there is a quantity of sesame above, to be multiplied by

  3 jin 7 fen 5 li. What is obtained?\\

  \textbf{Method:} Place the sesame quantity above; using 3 jin 7 fen 5 li as divisor,

  call from the second digit; when ten, place next to body; at head digit, break body; obtain.

}



\bilingual{%

  \textbf{問33.}

  今有小麥五碩九斛二升，每斛磨麵六斤一十四兩。問磨麵幾何？\\

  \textbf{答曰：}四百單七斤。\\

  \textbf{術曰：}列麥數於上，以六斤八分七釐半乘之，即得。

}{%

  \textbf{Problem 33.}

  Suppose there are 5 shuo 9 hu 2 sheng of wheat, each hu yielding

  6 jin 14 liang of flour. How much flour?\\

  \textbf{Answer:} 407 jin.\\

  \textbf{Method:} Place the wheat quantity above; multiply by 6 jin 8 fen 7 li

  and a half; obtain.\\

  (6 jin 14 liang = $6 + 14/16 = 6.875$ jin)

}



\bilingual{%

  \textbf{問34.}

  今有菉荳三十二碩七斛三升，每斛造粉五斤六兩。問造粉幾何？\\

  \textbf{答曰：}一千七百五十九斤三兩八錢。\\

  \textbf{術曰：}列荳數於上，以五斤六兩為法，從上次位呼之，即得。

}{%

  \textbf{Problem 34.}

  Suppose there are 32 shuo 7 hu 3 sheng of green beans, each hu producing

  5 jin 6 liang of bean flour. How much flour?\\

  \textbf{Answer:} 1,759 jin 3 liang 8 qian.\\

  \textbf{Method:} Place the bean quantity above; using 5 jin 6 liang as divisor,

  call from the second digit; obtain.

}



\bilingual{%

  \textbf{問35.}

  今有黃金一萬二千八百兩，每兩買地六畝二分五釐。問買地幾何？\\

  \textbf{答曰：}八萬畝。\\

  \textbf{術曰：}列金數於上，以地率從上次位呼之，即得。

}{%

  \textbf{Problem 35.}

  Suppose there are 12,800 liang of gold, each liang buying 6 mu 2 fen 5 li

  of land. How much land?\\

  \textbf{Answer:} 80,000 mu.\\

  \textbf{Method:} Place the gold quantity above; call the land rate from the

  second digit; obtain.\\

  (Calculation: $12{,}800 \times 6.25 = 80{,}000$)

}



\bilingual{%

  \textbf{問36.}

  今有田一百五十六頃二十五畝，每畝收稻五斛。問收稻幾何？\\

  \textbf{答曰：}九萬斛。\\

  \textbf{術曰：}列田數於上，以畝產從上次位呼之，即得。

}{%

  \textbf{Problem 36.}

  Suppose there are 156 qing 25 mu of field, each mu yielding 5 hu of rice.

  How much rice is harvested?\\

  \textbf{Answer:} 90,000 hu.\\

  \textbf{Method:} Place the field quantity above; call the per-mu yield from

  the second digit; obtain.

}



\bilingual{%

  \textbf{問37.}

  今有米五十三斛四升，直錢九貫八百七十九文。只有錢六十八貫二百六十文。

  問得米幾何？\\

  \textbf{答曰：}如術計之。\\

  \textbf{術曰：}列見錢為實，以米價為法，留頭乘之，即得。

}{%

  \textbf{Problem 37.}

  Suppose 53 hu 4 sheng of rice cost 9 guan 879 cash.

  If one has 68 guan 260 cash, how much rice can be bought?\\

  \textbf{Answer:} Calculate according to the method.\\

  \textbf{Method:} Place the available cash as dividend; use the rice price as

  divisor; keep-the-head multiply; obtain.

}



\bilingual{%

  \textbf{問38.}

  今有羅七百六十二匹，每匹價錢三貫二百三十文。問計錢幾何？\\

  \textbf{答曰：}如術計之。\\

  \textbf{術曰：}列羅數於上，以匹價從上次位呼之，即得。

}{%

  \textbf{Problem 38.}

  Suppose there are 762 bolts of gauze, each costing 3 guan 230 cash.

  How much cash in total?\\

  \textbf{Answer:} Calculate according to the method.\\

  \textbf{Method:} Place the gauze quantity above; call the bolt-price from

  the second digit; obtain.\\

  (Calculation: $762 \times 3{,}230 = 2{,}461{,}260$)

}



\longline





%% ========================================================================

%% Chapter 4: 身外減法門

%% ========================================================================

\bilingualsubsection{身外減法門}{Subtraction from the Body}

\label{sec:shenwaijian}



\bilingual{%

  \begin{rhymecn}

  減法根源必要知，即同求一一般推，\\

  呼如身下須當減，言十從身本位除。

  \end{rhymecn}

  身外減法者，除法之捷術也，與身外加法相反。

  其法以實為身，於身外逐位減之，呼如減下身位，言十從身本位除。

  凡十問。

}{%

  \begin{rhymecn}

  The root of the subtraction method you must know;\\

  It is like the ``seeking unity'' method applied;\\

  When ``as-is'' is called, subtract below the body;\\

  When ten is reached, divide from the body itself.

  \end{rhymecn}

  The ``subtraction from the body'' method is the inverse of the addition

  method: it performs division by subtracting from the dividend itself.

  ``As-is'' means subtract below the current digit; ``ten'' means subtract from

  the current digit itself. This is a streamlined division technique on the

  counting board.

}



\bilingual{%

  \textbf{問39.}

  今有錢四貫三百二十文，欲糴白豆，每斗價錢二十文。問得幾何？\\

  \textbf{答曰：}二百一十六斗。\\

  \textbf{術曰：}列錢數於上，為實，以斗價為法，從上身外減之，即得。

}{%

  \textbf{Problem 39.}

  Suppose there are 4 guan 320 cash, to buy white beans at 20 cash per dou.

  How much is obtained?\\

  \textbf{Answer:} 216 dou.\\

  \textbf{Method:} Place the cash above as dividend; use the dou-price as divisor;

  subtract from the body from above; obtain.\\

  (Calculation: $4{,}320 \div 20 = 216$)

}



\bilingual{%

  \textbf{問40.}

  今有錢四十三貫二百文，欲買細絲，每斤價錢三百文。問得幾何？\\

  \textbf{答曰：}一百四十四斤。\\

  \textbf{術曰：}列錢數於上，為實，以斤價為法，從上身外減之，即得。

}{%

  \textbf{Problem 40.}

  Suppose there are 43 guan 200 cash, to buy fine silk at 300 cash per jin.

  How much is obtained?\\

  \textbf{Answer:} 144 jin.\\

  \textbf{Method:} Place the cash above as dividend; use the jin-price as divisor;

  subtract from the body; obtain.\\

  (Calculation: $43{,}200 \div 300 = 144$)

}



\bilingual{%

  \textbf{問41.}

  今有錢一千四百一十六貫，欲買絹子，每匹價錢四貫文。問得幾何？\\

  \textbf{答曰：}三百五十四匹。\\

  \textbf{術曰：}列錢數於上，為實，以匹價為法，從上身外減之，即得。

}{%

  \textbf{Problem 41.}

  Suppose there are 1,416 guan, to buy silk cloth at 4 guan per bolt.

  How much is obtained?\\

  \textbf{Answer:} 354 bolts.\\

  \textbf{Method:} Place the cash above as dividend; use bolt-price as divisor;

  subtract from the body; obtain.\\

  (Calculation: $1{,}416 \div 4 = 354$)

}



\bilingual{%

  \textbf{問42.}

  今有銀二萬七千三百五十兩，欲為課銀，每錠五十兩。問為幾錠？\\

  \textbf{答曰：}五百四十七錠。\\

  \textbf{術曰：}列銀數於上，為實，以錠重為法，從上身外減之，即得。

}{%

  \textbf{Problem 42.}

  Suppose there are 27,350 liang of silver, to be made into tax ingots,

  each ingot = 50 liang. How many ingots?\\

  \textbf{Answer:} 547 ingots.\\

  \textbf{Method:} Place the silver above as dividend; use ingot weight as divisor;

  subtract from the body; obtain.\\

  (Calculation: $27{,}350 \div 50 = 547$)

}



\bilingual{%

  \textbf{問43.}

  今有錢四千四百一十六貫，欲買大羅，每匹價錢六貫文。問得幾何？\\

  \textbf{答曰：}七百三十六匹。\\

  \textbf{術曰：}列錢數於上，為實，以匹價為法，從上身外減之，即得。

}{%

  \textbf{Problem 43.}

  Suppose there are 4,416 guan, to buy damask at 6 guan per bolt.

  How much is obtained?\\

  \textbf{Answer:} 736 bolts.\\

  \textbf{Method:} Place the cash above as dividend; use bolt-price as divisor;

  subtract from the body; obtain.\\

  (Calculation: $4{,}416 \div 6 = 736$)

}



\bilingual{%

  \textbf{問44.}

  今有錢六百二十四貫四百文，欲買甘草，每秤價錢七百文。問得幾何？\\

  \textbf{答曰：}八百九十二秤。\\

  \textbf{術曰：}列錢數於上，為實，以秤價為法，從上身外減之，即得。

}{%

  \textbf{Problem 44.}

  Suppose there are 624 guan 400 cash, to buy licorice at 700 cash per cheng.

  How much is obtained?\\

  \textbf{Answer:} 892 cheng.\\

  \textbf{Method:} Place the cash above as dividend; use cheng-price as divisor;

  subtract from the body; obtain.\\

  (Calculation: $624{,}400 \div 700 = 892$)

}



\bilingual{%

  \textbf{問45.}

  今有錢五十貫七百二十文，欲買細布，每尺價錢八十文。問得幾何？\\

  \textbf{答曰：}六百三十四尺。\\

  \textbf{術曰：}列錢數於上，為實，以尺價為法，從上身外減之，即得。

}{%

  \textbf{Problem 45.}

  Suppose there are 50 guan 720 cash, to buy fine cloth at 80 cash per chi.

  How much is obtained?\\

  \textbf{Answer:} 634 chi.\\

  \textbf{Method:} Place the cash above as dividend; use chi-price as divisor;

  subtract from the body; obtain.\\

  (Calculation: $50{,}720 \div 80 = 634$)

}



\bilingual{%

  \textbf{問46.}

  今有錢三萬八千二百五十貫，欲買馬，每匹價錢九十貫。問得幾何？\\

  \textbf{答曰：}四百二十五匹。\\

  \textbf{術曰：}列錢數於上，為實，以匹價為法，從上身外減之，即得。

}{%

  \textbf{Problem 46.}

  Suppose there are 38,250 guan, to buy horses at 90 guan each.

  How many horses?\\

  \textbf{Answer:} 425 horses.\\

  \textbf{Method:} Place the cash above as dividend; use horse-price as divisor;

  subtract from the body; obtain.\\

  (Calculation: $38{,}250 \div 90 = 425$)

}



\bilingual{%

  \textbf{問47.}

  今有油六碩八斗四升，每三斤七分五釐直錢一百文。問直錢幾何？\\

  \textbf{答曰：}如術計之。\\

  \textbf{術曰：}列油數為實，以三斤七分五釐為法，身外減之，即得。

}{%

  \textbf{Problem 47.}

  Suppose there are 6 shuo 8 dou 4 sheng of oil, and 3 jin 7 fen 5 li costs

  100 cash. How much cash?\\

  \textbf{Answer:} Calculate according to the method.\\

  \textbf{Method:} Place the oil quantity as dividend; use 3 jin 7 fen 5 li as

  divisor; subtract from the body; obtain.

}



\bilingual{%

  \textbf{問48.}

  今有麵四百單七斤，每六斤八分七釐半直錢一百文。問直錢幾何？\\

  \textbf{答曰：}如術計之。\\

  \textbf{術曰：}列麵數為實，以六斤八分七釐半為法，身外減之，即得。

}{%

  \textbf{Problem 48.}

  Suppose there are 407 jin of flour, and 6 jin 8 fen 7.5 li costs 100 cash.

  How much cash?\\

  \textbf{Answer:} Calculate according to the method.\\

  \textbf{Method:} Place the flour quantity as dividend; use 6 jin 8 fen 7.5 li as

  divisor; subtract from the body; obtain.

}



\longline





%% ========================================================================

%% Chapter 5: 九歸除法門

%% ========================================================================

\bilingualsubsection{九歸除法門}{Division by the Nine Returns}

\label{sec:jiugui}



\bilingual{%

  \begin{rhymecn}

  實少法多從法歸，實多滿法進前居，\\

  常存除數專心記，法實相停九十餘。\\

  但遇無除還頭位，然將釋九數呼除，\\

  流傳故泄真消息，求一穿韜總不如。

  \end{rhymecn}

  九歸除法者，以九歸訣為除法之要術也。

  實少於法，則從法歸之；實多於法，則進位而居之。

  常記除數於心，法實相等則餘九十餘。

  凡二十九問。

}{%

  \begin{rhymecn}

  When the dividend is small and divisor large, follow the return method;\\

  When the dividend fills the divisor, advance to the next position;\\

  Keep the divisor always in mind;\\

  When divisor and dividend match, the remainder is in the nineties.\\

  When a digit cannot be divided, return to the head position;\\

  Then use the nine-explanations method to call the division;\\

  This tradition reveals the true secret;\\

  The ``seeking unity'' method cannot compare.

  \end{rhymecn}

  The ``nine returns'' division uses the division rhymes given in the

  General Principles. When the dividend is smaller than the divisor, use the

  ``return'' method; when the dividend is larger, advance the quotient. The

  practitioner must memorize the divisor and keep it in mind. When the division

  is exact, the remainder is in the nineties. If a digit cannot be divided,

  return to the head position and use the ``nine explanations'' method. This

  method, though not the orthodox approach (商除), is easier for beginners.

}



\bilingual{%

  \textbf{問49.}

  今有錢四貫三百二十文，欲糴白豆，每斗價錢二十文。問得幾何？\\

  \textbf{答曰：}二百一十六斗。\\

  \textbf{術曰：}列錢物於上為實，各以價錢為法，從上身外減之，即得。

}{%

  \textbf{Problem 49.}

  Suppose there are 4 guan 320 cash, to buy white beans at 20 cash per dou.

  How much is obtained?\\

  \textbf{Answer:} 216 dou.\\

  \textbf{Method:} Place cash and goods above as dividend; use the price as divisor;

  subtract from the body; obtain.

}



\bilingual{%

  \textbf{問50.}

  今有錢四十三貫二百文，欲買細絲，每斤價錢三百文。問得幾何？\\

  \textbf{答曰：}一百四十四斤。\\

  \textbf{術曰：}列錢數為實，以斤價為法，九歸除之，即得。

}{%

  \textbf{Problem 50.}

  Suppose there are 43 guan 200 cash, to buy fine silk at 300 cash per jin.

  How much is obtained?\\

  \textbf{Answer:} 144 jin.\\

  \textbf{Method:} Place cash as dividend; use jin-price as divisor; divide by

  nine-returns; obtain.

}



\bilingual{%

  \textbf{問51.}

  今有錢一千四百一十六貫，欲買絹子，每匹價錢四貫文。問得幾何？\\

  \textbf{答曰：}三百五十四匹。\\

  \textbf{術曰：}列錢數為實，以匹價為法，九歸除之，即得。

}{%

  \textbf{Problem 51.}

  Suppose there are 1,416 guan, to buy silk at 4 guan per bolt.

  How much is obtained?\\

  \textbf{Answer:} 354 bolts.\\

  \textbf{Method:} Place cash as dividend; use bolt-price as divisor; divide by

  nine-returns; obtain.

}



\bilingual{%

  \textbf{問52.}

  今有銀二萬七千三百五十兩，欲為課銀，每錠五十兩。問為幾錠？\\

  \textbf{答曰：}五百四十七錠。\\

  \textbf{術曰：}列銀數為實，以錠重為法，九歸除之，即得。

}{%

  \textbf{Problem 52.}

  Suppose there are 27,350 liang of silver, to make tax ingots of 50 liang each.

  How many ingots?\\

  \textbf{Answer:} 547 ingots.\\

  \textbf{Method:} Place silver as dividend; use ingot weight as divisor;

  divide by nine-returns; obtain.

}



\bilingual{%

  \textbf{問53.}

  今有錢四千四百一十六貫，欲買大羅，每匹價錢六貫文。問得幾何？\\

  \textbf{答曰：}七百三十六匹。\\

  \textbf{術曰：}列錢數為實，以匹價為法，九歸除之，即得。

}{%

  \textbf{Problem 53.}

  Suppose there are 4,416 guan, to buy damask at 6 guan per bolt.

  How much is obtained?\\

  \textbf{Answer:} 736 bolts.\\

  \textbf{Method:} Place cash as dividend; use bolt-price as divisor; divide by

  nine-returns; obtain.

}



\bilingual{%

  \textbf{問54.}

  今有錢六百二十四貫四百文，欲買甘草，每秤價錢七百文。問得幾何？\\

  \textbf{答曰：}八百九十二秤。\\

  \textbf{術曰：}列錢數為實，以秤價為法，九歸除之，即得。

}{%

  \textbf{Problem 54.}

  Suppose there are 624 guan 400 cash, to buy licorice at 700 cash per cheng.

  How much is obtained?\\

  \textbf{Answer:} 892 cheng.\\

  \textbf{Method:} Place cash as dividend; use cheng-price as divisor; divide by

  nine-returns; obtain.

}



\bilingual{%

  \textbf{問55.}

  今有錢五十貫七百二十文，欲買細布，每尺價錢八十文。問得幾何？\\

  \textbf{答曰：}六百三十四尺。\\

  \textbf{術曰：}列錢數為實，以尺價為法，九歸除之，即得。

}{%

  \textbf{Problem 55.}

  Suppose there are 50 guan 720 cash, to buy fine cloth at 80 cash per chi.

  How much is obtained?\\

  \textbf{Answer:} 634 chi.\\

  \textbf{Method:} Place cash as dividend; use chi-price as divisor; divide by

  nine-returns; obtain.

}



\bilingual{%

  \textbf{問56.}

  今有木幾何，直錢三千二百六十株。問每株價錢幾何？\\

  \textbf{答曰：}三貫二百六十文。\\

  \textbf{術曰：}列錢物於上為實，各以價錢為法，從上身外減之，即得。

}{%

  \textbf{Problem 56.}

  Suppose some quantity of timber costs 3,260 trees' worth.

  What is the price per tree?\\

  \textbf{Answer:} 3 guan 260 cash.\\

  \textbf{Method:} Place cash and goods above as dividend; use price as divisor;

  subtract from the body; obtain.

}



\bilingual{%

  \textbf{問57.}

  今有絹幾何，每匹四貫，共錢一千四百一十六貫。問匹數幾何？\\

  \textbf{答曰：}三百五十四匹。\\

  \textbf{術曰：}列錢數為實，以匹價為法，九歸除之，即得。

}{%

  \textbf{Problem 57.}

  Suppose there is some silk, 4 guan per bolt, total 1,416 guan.

  How many bolts?\\

  \textbf{Answer:} 354 bolts.\\

  \textbf{Method:} Place cash as dividend; use bolt-price as divisor; divide by

  nine-returns; obtain.

}



\bilingual{%

  \textbf{問58.}

  今有銀幾何，每錠五十兩，共二萬七千三百五十兩。問錠數幾何？\\

  \textbf{答曰：}五百四十七錠。\\

  \textbf{術曰：}列銀數為實，以錠重為法，九歸除之，即得。

}{%

  \textbf{Problem 58.}

  Suppose there is some silver, 50 liang per ingot, total 27,350 liang.

  How many ingots?\\

  \textbf{Answer:} 547 ingots.\\

  \textbf{Method:} Place silver as dividend; use ingot weight as divisor;

  divide by nine-returns; obtain.

}



\bilingual{%

  \textbf{問59.}

  今有甘草幾何，每秤七百文，共錢六百二十四貫四百文。問秤數幾何？\\

  \textbf{答曰：}八百九十二秤。\\

  \textbf{術曰：}列錢數為實，以秤價為法，九歸除之，即得。

}{%

  \textbf{Problem 59.}

  Suppose there is some licorice, 700 cash per cheng, total 624 guan 400 cash.

  How many cheng?\\

  \textbf{Answer:} 892 cheng.\\

  \textbf{Method:} Place cash as dividend; use cheng-price as divisor; divide by

  nine-returns; obtain.

}



\bilingual{%

  \textbf{問60.}

  今有麻布幾何，每匹一百九十四文，共錢一百四貫三百七十二文。問匹數幾何？\\

  \textbf{答曰：}五百三十八匹。\\

  \textbf{術曰：}列錢數為實，以匹價為法，九歸除之，即得。

}{%

  \textbf{Problem 60.}

  Suppose there is some hemp cloth, 194 cash per bolt, total 104 guan 372 cash.

  How many bolts?\\

  \textbf{Answer:} 538 bolts.\\

  \textbf{Method:} Place cash as dividend; use bolt-price as divisor; divide by

  nine-returns; obtain.

}



\bilingual{%

  \textbf{問61.}

  今有細布幾何，每尺八十文，共錢五十貫七百二十文。問尺數幾何？\\

  \textbf{答曰：}六百三十四尺。\\

  \textbf{術曰：}列錢數為實，以尺價為法，九歸除之，即得。

}{%

  \textbf{Problem 61.}

  Suppose there is some fine cloth, 80 cash per chi, total 50 guan 720 cash.

  How many chi?\\

  \textbf{Answer:} 634 chi.\\

  \textbf{Method:} Place cash as dividend; use chi-price as divisor; divide by

  nine-returns; obtain.

}



\bilingual{%

  \textbf{問62.}

  今有松木幾何，每株一貫七百五文，共錢五千五百五十八貫三百文。問株數幾何？\\

  \textbf{答曰：}三千二百六十株。\\

  \textbf{術曰：}列錢數為實，以株價為法，九歸除之，即得。

}{%

  \textbf{Problem 62.}

  Suppose there are some pine trees, 1 guan 705 cash each, total 5,558 guan 300 cash.

  How many trees?\\

  \textbf{Answer:} 3,260 trees.\\

  \textbf{Method:} Place cash as dividend; use tree-price as divisor; divide by

  nine-returns; obtain.

}



\bilingual{%

  \textbf{問63.}

  今有白豆幾何，每斗二十文，共錢四貫三百二十文。問斗數幾何？\\

  \textbf{答曰：}二百一十六斗。\\

  \textbf{術曰：}列錢數為實，以斗價為法，九歸除之，即得。

}{%

  \textbf{Problem 63.}

  Suppose there is some white beans, 20 cash per dou, total 4 guan 320 cash.

  How many dou?\\

  \textbf{Answer:} 216 dou.\\

  \textbf{Method:} Place cash as dividend; use dou-price as divisor; divide by

  nine-returns; obtain.

}



\bilingual{%

  \textbf{問64.}

  今有絲幾何，每斤三百文，共錢四十三貫二百文。問斤數幾何？\\

  \textbf{答曰：}一百四十四斤。\\

  \textbf{術曰：}列錢數為實，以斤價為法，九歸除之，即得。

}{%

  \textbf{Problem 64.}

  Suppose there is some silk, 300 cash per jin, total 43 guan 200 cash.

  How many jin?\\

  \textbf{Answer:} 144 jin.\\

  \textbf{Method:} Place cash as dividend; use jin-price as divisor; divide by

  nine-returns; obtain.

}



\bilingual{%

  \textbf{問65.}

  今有馬幾何，每匹九十貫，共錢三萬八千二百五十貫。問匹數幾何？\\

  \textbf{答曰：}四百二十五匹。\\

  \textbf{術曰：}列錢數為實，以匹價為法，九歸除之，即得。

}{%

  \textbf{Problem 65.}

  Suppose there are some horses, 90 guan each, total 38,250 guan.

  How many horses?\\

  \textbf{Answer:} 425 horses.\\

  \textbf{Method:} Place cash as dividend; use horse-price as divisor; divide by

  nine-returns; obtain.

}



\bilingual{%

  \textbf{問66.}

  今有羅幾何，每匹價錢三貫二百三十文，共錢一千七百一十一貫四百文。問匹數幾何？\\

  \textbf{答曰：}如術計之。\\

  \textbf{術曰：}列錢數為實，以匹價為法，九歸除之，即得。

}{%

  \textbf{Problem 66.}

  Suppose there is some gauze, 3 guan 230 cash per bolt,

  total 1,711 guan 400 cash. How many bolts?\\

  \textbf{Answer:} Calculate according to the method.\\

  \textbf{Method:} Place cash as dividend; use bolt-price as divisor; divide by

  nine-returns; obtain.

}



\bilingual{%

  \textbf{問67.}

  今有大羅幾何，每匹六貫，共錢四千四百一十六貫。問匹數幾何？\\

  \textbf{答曰：}七百三十六匹。\\

  \textbf{術曰：}列錢數為實，以匹價為法，九歸除之，即得。

}{%

  \textbf{Problem 67.}

  Suppose there is some damask, 6 guan per bolt, total 4,416 guan.

  How many bolts?\\

  \textbf{Answer:} 736 bolts.\\

  \textbf{Method:} Place cash as dividend; use bolt-price as divisor; divide by

  nine-returns; obtain.

}



\bilingual{%

  \textbf{問68.}

  今有苧絲幾何，每尺一百九十文，共錢六十五貫五百五十文。問尺數幾何？\\

  \textbf{答曰：}三百四十五尺。\\

  \textbf{術曰：}列錢數為實，以尺價為法，九歸除之，即得。

}{%

  \textbf{Problem 68.}

  Suppose there is some ramie silk, 190 cash per chi, total 65 guan 550 cash.

  How many chi?\\

  \textbf{Answer:} 345 chi.\\

  \textbf{Method:} Place cash as dividend; use chi-price as divisor; divide by

  nine-returns; obtain.

}



\bilingual{%

  \textbf{問69.}

  今有米幾何，每斗一百一十文，共錢七貫五百二十四文。問斗數幾何？\\

  \textbf{答曰：}六碩八斗四升。\\

  \textbf{術曰：}列錢數為實，以斗價為法，九歸除之，即得。

}{%

  \textbf{Problem 69.}

  Suppose there is some rice, 110 cash per dou, total 7 guan 524 cash.

  How many dou?\\

  \textbf{Answer:} 6 shuo 8 dou 4 sheng.\\

  \textbf{Method:} Place cash as dividend; use dou-price as divisor; divide by

  nine-returns; obtain.\\

  (Calculation: $7{,}524 \div 110 = 68.4$ dou = 6 shuo 8 dou 4 sheng)

}



\bilingual{%

  \textbf{問70.}

  今有鹽幾何，每袋一十三貫，共錢一萬一千三百四十九貫。問袋數幾何？\\

  \textbf{答曰：}八百七十三袋。\\

  \textbf{術曰：}列錢數為實，以袋價為法，九歸除之，即得。

}{%

  \textbf{Problem 70.}

  Suppose there is some salt, 13 guan per bag, total 11,349 guan.

  How many bags?\\

  \textbf{Answer:} 873 bags.\\

  \textbf{Method:} Place cash as dividend; use bag-price as divisor; divide by

  nine-returns; obtain.

}



\bilingual{%

  \textbf{問71.}

  今有羊幾何，每只四貫，共錢一千四百一十六貫。問只數幾何？\\

  \textbf{答曰：}三百五十四只。\\

  \textbf{術曰：}列錢數為實，以只價為法，九歸除之，即得。

}{%

  \textbf{Problem 71.}

  Suppose there are some sheep, 4 guan each, total 1,416 guan.

  How many sheep?\\

  \textbf{Answer:} 354 sheep.\\

  \textbf{Method:} Place cash as dividend; use sheep-price as divisor; divide by

  nine-returns; obtain.

}



\bilingual{%

  \textbf{問72.}

  今有油六碩八斗四升，每三斤七分五釐直錢一百文。問直錢幾何？\\

  \textbf{答曰：}如術計之。\\

  \textbf{術曰：}列油數為實，以三斤七分五釐為法，九歸除之，即得。

}{%

  \textbf{Problem 72.}

  Suppose there are 6 shuo 8 dou 4 sheng of oil; 3 jin 7 fen 5 li costs 100 cash.

  How much cash?\\

  \textbf{Answer:} Calculate according to the method.\\

  \textbf{Method:} Place the oil quantity as dividend; use 3 jin 7 fen 5 li as

  divisor; divide by nine-returns; obtain.

}



\bilingual{%

  \textbf{問73.}

  今有麵四百單七斤，每六斤八分七釐半直錢一百文。問直錢幾何？\\

  \textbf{答曰：}如術計之。\\

  \textbf{術曰：}列麵數為實，以六斤八分七釐半為法，九歸除之，即得。

}{%

  \textbf{Problem 73.}

  Suppose there are 407 jin of flour; 6 jin 8 fen 7.5 li costs 100 cash.

  How much cash?\\

  \textbf{Answer:} Calculate according to the method.\\

  \textbf{Method:} Place the flour quantity as dividend; use 6 jin 8 fen 7.5 li as

  divisor; divide by nine-returns; obtain.

}



\bilingual{%

  \textbf{問74.}

  今有羅三十四尺六寸，每尺一百二十文。問計錢幾何？\\

  \textbf{答曰：}四貫一百五十二文。\\

  \textbf{術曰：}列羅數於上，以尺價為法，九歸除之，即得。

}{%

  \textbf{Problem 74.}

  Suppose there are 34 chi 6 cun of gauze, each chi costing 120 cash.

  How much cash in total?\\

  \textbf{Answer:} 4 guan 152 cash.\\

  \textbf{Method:} Place the gauze quantity above; use chi-price as divisor;

  divide by nine-returns; obtain.

}



\bilingual{%

  \textbf{問75.}

  今有胡椒六十三斤四兩，每斤三百八十文。問計錢幾何？\\

  \textbf{答曰：}二十四貫三十五文。\\

  \textbf{術曰：}列椒數於上，以斤價為法，九歸除之，即得。

}{%

  \textbf{Problem 75.}

  Suppose there are 63 jin 4 liang of pepper, each jin costing 380 cash.

  How much cash in total?\\

  \textbf{Answer:} 24 guan 35 cash.\\

  \textbf{Method:} Place the pepper quantity above; use jin-price as divisor;

  divide by nine-returns; obtain.

}



\bilingual{%

  \textbf{問76.}

  今有馬幾何，每匹價錢九十貫。問計錢幾何？\\

  \textbf{答曰：}如術計之。\\

  \textbf{術曰：}列馬數於上，以匹價為法，九歸除之，即得。

}{%

  \textbf{Problem 76.}

  Suppose there are some horses, each costing 90 guan.

  How much cash in total?\\

  \textbf{Answer:} Calculate according to the method.\\

  \textbf{Method:} Place the horse quantity above; use horse-price as divisor;

  divide by nine-returns; obtain.

}



\bilingual{%

  \textbf{問77.}

  今有綿三百四十五斤，欲作綿被，每床用綿七斤。問作幾床？\\

  \textbf{答曰：}四十九床，餘二斤。\\

  \textbf{術曰：}列綿數為實，以每床用綿為法，九歸除之，即得。

}{%

  \textbf{Problem 77.}

  Suppose there are 345 jin of cotton, to make cotton quilts,

  each quilt using 7 jin. How many quilts?\\

  \textbf{Answer:} 49 quilts, remainder 2 jin.\\

  \textbf{Method:} Place the cotton quantity as dividend; use cotton per quilt as

  divisor; divide by nine-returns; obtain.\\

  (Calculation: $345 \div 7 = 49$ remainder $2$)

}



\longline





%% ========================================================================

%% Chapter 6: 異乘同除門

%% ========================================================================

\bilingualsubsection{異乘同除門}{Different Multiplication, Same Division}

\label{sec:yicheng}



\bilingual{%

  異乘同除者，比例交易之法也。以異名之數相乘，以同名之數相除，

  求其相當之率。凡八問。

}{%

  This chapter covers problems of proportional exchange: multiplying by one

  rate and dividing by another. The problems involve currency conversion,

  unit conversion, and commodity exchange.

}



\bilingual{%

  \textbf{問78.}

  今有錢九貫八百七十九文，糴米五碩三斛四升。只有錢六十八貫二百六十文。

  問得米幾何？\\

  \textbf{答曰：}如術計之。\\

  \textbf{術曰：}列只有米數，以乘今有錢為實，以原錢為法，除之，即得。

}{%

  \textbf{Problem 78.}

  Suppose 9 guan 879 cash buys 5 shuo 3 hu 4 sheng of rice.

  If one has 68 guan 260 cash, how much rice is obtained?\\

  \textbf{Answer:} Calculate according to the method.\\

  \textbf{Method:} Place the rice quantity, multiply by the current cash as dividend;

  use the original cash as divisor; divide; obtain.

}



\bilingual{%

  \textbf{問79.}

  今有絲七匹，通尺內子得二千三十一尺，以乘上位得六百三十七萬九千三百七十一。

  又以實以斤銖法三百八十四銖除之，得三萬六千五百四十二銖，為實。\\

  \textbf{答曰：}九千七百六十五斤一十兩四銖。\\

  \textbf{術曰：}列七匹通尺內子，以乘上位，又以斤銖法除之，即得。

}{%

  \textbf{Problem 79.}

  Suppose there are 7 bolts of silk; converting to chi gives 2,031 chi;

  multiplying gives 6,379,371. Then dividing by 384 zhu (the jin-zhu conversion)

  gives 36,542 zhu as dividend.\\

  \textbf{Answer:} 9,765 jin 10 liang 4 zhu.\\

  \textbf{Method:} Convert 7 bolts to chi, multiply, then divide by the jin-zhu

  conversion factor; obtain.

}



\bilingual{%

  \textbf{問80.}

  今有絲八斤二兩二十一銖，織錦七匹六尺五寸。只有絲九十五斤三兩六銖。

  問織錦幾何？\\

  \textbf{答曰：}如術計之。\\

  \textbf{術曰：}（匹法同前）列只有絲數，以乘錦數為實，以原絲為法，除之，即得。

}{%

  \textbf{Problem 80.}

  Suppose 8 jin 2 liang 21 zhu of silk weaves 7 bolts 6 chi 5 cun of brocade.

  If one has 95 jin 3 liang 6 zhu of silk, how much brocade can be woven?\\

  \textbf{Answer:} Calculate according to the method.\\

  \textbf{Method:} (Bolt conversion same as before) Place the available silk,

  multiply by brocade quantity as dividend, use original silk as divisor; divide; obtain.

}



\bilingual{%

  \textbf{問81.}

  今有錢五貫八百三十文，買麩八斗五升。只有錢三十七貫七百六十文。

  問買麩幾何？\\

  \textbf{答曰：}五碩五斗二升。\\

  \textbf{術曰：}列只有錢數，以乘原麩為實，以原錢為法，除之，即得。

}{%

  \textbf{Problem 81.}

  Suppose 5 guan 830 cash buys 8 dou 5 sheng of bran.

  If one has 37 guan 760 cash, how much bran can be bought?\\

  \textbf{Answer:} 5 shuo 5 dou 2 sheng.\\

  \textbf{Method:} Place the available cash, multiply by original bran as dividend;

  use original cash as divisor; divide; obtain.

}



\bilingual{%

  \textbf{問82.}

  今有錢一貫一百七十五文，買麥五斗四升。只有錢二十三貫五百文。

  問買麥幾何？\\

  \textbf{答曰：}十碩八斗。\\

  \textbf{術曰：}列只有錢數，以乘原麥為實，以原錢為法，除之，即得。

}{%

  \textbf{Problem 82.}

  Suppose 1 guan 175 cash buys 5 dou 4 sheng of wheat.

  If one has 23 guan 500 cash, how much wheat can be bought?\\

  \textbf{Answer:} 10 shuo 8 dou.\\

  \textbf{Method:} Place available cash, multiply by original wheat as dividend;

  use original cash as divisor; divide; obtain.\\

  (Calculation: $\dfrac{23{,}500 \times 5.4}{1{,}175} = 108$ dou = 10 shuo 8 dou)

}



\bilingual{%

  \textbf{問83.}

  今有錢二貫六百文，買麵三斤四兩。只有錢四十七貫四百五十文。

  問買麵幾何？\\

  \textbf{答曰：}六十三斤四兩。\\

  \textbf{術曰：}列只有錢數，以乘原麵為實，以原錢為法，除之，即得。

}{%

  \textbf{Problem 83.}

  Suppose 2 guan 600 cash buys 3 jin 4 liang of flour.

  If one has 47 guan 450 cash, how much flour can be bought?\\

  \textbf{Answer:} 63 jin 4 liang.\\

  \textbf{Method:} Place available cash, multiply by original flour as dividend;

  use original cash as divisor; divide; obtain.

}



\bilingual{%

  \textbf{問84.}

  今有錢一貫二百文，買鹽二碩四斗。只有錢三十二貫三百文。

  問買鹽幾何？\\

  \textbf{答曰：}六十四碩八斗。\\

  \textbf{術曰：}列只有錢數，以乘原鹽為實，以原錢為法，除之，即得。

}{%

  \textbf{Problem 84.}

  Suppose 1 guan 200 cash buys 2 shuo 4 dou of salt.

  If one has 32 guan 300 cash, how much salt can be bought?\\

  \textbf{Answer:} 64 shuo 8 dou.\\

  \textbf{Method:} Place available cash, multiply by original salt as dividend;

  use original cash as divisor; divide; obtain.\\

  (Calculation: $\dfrac{32{,}300 \times 24}{1{,}200} = 648$ dou = 64 shuo 8 dou)

}



\bilingual{%

  \textbf{問85.}

  今有銀二十四銖，直錢三貫六百文。只有銀五十九兩八銖。

  問直錢幾何？\\

  \textbf{答曰：}八十九貫七百文。\\

  \textbf{術曰：}列只有銀數，以乘原直錢為實，以原銀為法，除之，即得。

}{%

  \textbf{Problem 85.}

  Suppose 24 zhu of silver is worth 3 guan 600 cash.

  If one has 59 liang 8 zhu of silver, what is it worth?\\

  \textbf{Answer:} 89 guan 700 cash.\\

  \textbf{Method:} Place the available silver, multiply by original cash value as

  dividend; use original silver as divisor; divide; obtain.

}



\longline



%% ========================================================================

%% Chapter 7: 庫務解稅門

%% ========================================================================

\bilingualsubsection{庫務解稅門}{Government Treasuries and Taxation}

\label{sec:kuwu}



\bilingual{%

  庫務解稅者，官司會計之法也。凡課銀、鈔法、子母相生、

  合金之色、工墨之費，皆出於此。凡十一問。

}{%

  This chapter deals with problems of government finance: calculating

  taxes, interest, alloy composition of currency, and distribution of resources.

  These problems reflect the fiscal administration of the Yuan dynasty.

}



\bilingual{%

  \textbf{問86.}

  今有本銀四千五百六兩，經三箇月一十二日，月利銀三百六兩。

  問本利共幾何？\\

  \textbf{答曰：}本銀四千五十兩，利銀三百六兩。\\

  \textbf{術曰：}置本銀以月數乘之，又以日數乘之，以三十日除之，得利息，

  加入本銀，即得。

}{%

  \textbf{Problem 86.}

  Suppose there is principal silver of 4,506 liang, over 3 months and 12 days,

  with monthly interest of 306 liang. What is the total of principal and interest?\\

  \textbf{Answer:} Principal 4,050 liang, interest 306 liang.\\

  \textbf{Method:} Place the principal; multiply by months, then multiply by days;

  divide by 30 days to get interest; add to principal; obtain.

}



\bilingual{%

  \textbf{問87.}

  今有本銀四千五十兩，月利一分二釐。問三箇月一十二日，利銀幾何？\\

  \textbf{答曰：}一百四十五兩八錢。\\

  \textbf{術曰：}置本銀以月利乘之，又以日數乘之，以三十日除之，即得。

}{%

  \textbf{Problem 87.}

  Suppose there is principal silver of 4,050 liang, monthly interest 1 fen 2 li

  (1.2\%). Over 3 months and 12 days, what is the interest?\\

  \textbf{Answer:} 145 liang 8 qian.\\

  \textbf{Method:} Place principal; multiply by monthly interest, then multiply by

  days; divide by 30 days; obtain.

}



\bilingual{%

  \textbf{問88.}

  今有課銀三萬六千兩，每五十兩為一錠。問為幾錠？\\

  \textbf{答曰：}七百二十錠。\\

  \textbf{術曰：}置課銀以錠重除之，即得。

}{%

  \textbf{Problem 88.}

  Suppose there are 36,000 liang of tax silver; each ingot = 50 liang.

  How many ingots?\\

  \textbf{Answer:} 720 ingots.\\

  \textbf{Method:} Place the tax silver; divide by ingot weight; obtain.\\

  (Calculation: $36{,}000 \div 50 = 720$)

}



\bilingual{%

  \textbf{問89.}

  今有鈔二千四百貫，欲換銀，每銀一兩直鈔八貫。問得銀幾何？\\

  \textbf{答曰：}三百兩。\\

  \textbf{術曰：}置鈔數為實，以銀價為法，除之，即得。

}{%

  \textbf{Problem 89.}

  Suppose there are 2,400 guan of paper money (鈔), to exchange for silver;

  each liang of silver is worth 8 guan of cash. How much silver?\\

  \textbf{Answer:} 300 liang.\\

  \textbf{Method:} Place the cash quantity as dividend; use silver price as divisor;

  divide; obtain.\\

  (Calculation: $2{,}400 \div 8 = 300$)

}



\bilingual{%

  \textbf{問90.}

  今有銀三百兩，欲換鈔，每兩直鈔八貫。問得鈔幾何？\\

  \textbf{答曰：}二千四百貫。\\

  \textbf{術曰：}置銀數為實，以鈔價為法，乘之，即得。

}{%

  \textbf{Problem 90.}

  Suppose there are 300 liang of silver, to exchange for paper money;

  each liang is worth 8 guan. How much cash?\\

  \textbf{Answer:} 2,400 guan.\\

  \textbf{Method:} Place the silver quantity as dividend; use cash price as divisor;

  multiply; obtain.\\

  (Calculation: $300 \times 8 = 2{,}400$)

}



\bilingual{%

  \textbf{問91.}

  今有金一十二兩五錢，內減五十兩，餘即入銀，合同為法。

  實如法而一，得本銀，反減共還銀數，餘即利銀也。\\

  \textbf{答曰：}一十二兩五錢。\\

  \textbf{術曰：}列金數為實，以八分為法，除之，得本銀，反減共還銀數，

  餘即利銀也。

}{%

  \textbf{Problem 91.}

  Suppose there is gold of 12 liang 5 qian; subtract 50 liang from it,

  the remainder enters into silver; combine as divisor.

  Divide the dividend by the divisor; obtain the principal silver.

  Subtract from the total returned silver; the remainder is the interest.\\

  \textbf{Answer:} 12 liang 5 qian.\\

  \textbf{Method:} Place the gold quantity as dividend; use 8 fen as divisor;

  divide to get principal silver; subtract from total returned silver;

  the remainder is interest.

}



\bilingual{%

  \textbf{問92.}

  今有銀一十二兩五錢，內減五十兩，餘即入銀，合同為法。

  實如法而一，得六兩，問為顏色分數幾何？\\

  \textbf{答曰：}八分色。\\

  \textbf{術曰：}列金數為實，以本銀為法，除之，即得顏色分數。

}{%

  \textbf{Problem 92.}

  Suppose there is silver of 12 liang 5 qian; subtract 50 liang,

  the remainder enters silver; combine as divisor.

  Dividing gives 6 liang; what is the alloy purity?\\

  \textbf{Answer:} 8 fen purity (80\% pure).\\

  \textbf{Method:} Place the gold quantity as dividend; use principal silver as

  divisor; divide; obtain the purity fraction.

}



\bilingual{%

  \textbf{問93.}

  今有絲六十二斤半，換金一十兩，內八分半金五兩七分半金五兩。

  問二色金各得幾何？\\

  \textbf{術曰：}列絲數為實，以金價為法，除之，即得各金數。

}{%

  \textbf{Problem 93.}

  Suppose 62.5 jin of silk is exchanged for 10 liang of gold:

  5 liang of 8.5-fen gold and 5 liang of 7.5-fen gold.

  How much of each type of gold?\\

  \textbf{Method:} Place the silk quantity as dividend; use gold price as divisor;

  divide; obtain each gold quantity.

}



\bilingual{%

  \textbf{問94.}

  今有銀三千五百六兩，直錢三十五貫六百文。問每兩直錢幾何？\\

  \textbf{答曰：}一十文。\\

  \textbf{術曰：}置銀數為實，以直錢為法，除之，即得。

}{%

  \textbf{Problem 94.}

  Suppose 3,506 liang of silver is worth 35 guan 600 cash.

  How much cash per liang?\\

  \textbf{Answer:} 10 cash.\\

  \textbf{Method:} Place the silver quantity as dividend; use cash value as divisor;

  divide; obtain.

}



\bilingual{%

  \textbf{問95.}

  今有鈔一百六十五貫五百文，買絲三十七斤八兩。問每斤價鈔幾何？\\

  \textbf{答曰：}四貫四百文。\\

  \textbf{術曰：}置鈔數為實，以絲數為法，除之，即得。

}{%

  \textbf{Problem 95.}

  Suppose 165 guan 500 cash buys 37 jin 8 liang of silk.

  What is the price per jin?\\

  \textbf{Answer:} 4 guan 400 cash per jin.\\

  \textbf{Method:} Place the cash quantity as dividend; use silk quantity as divisor;

  divide; obtain.

}



\bilingual{%

  \textbf{問96.}

  今有課鈔一千貫，納官每百貫收工墨錢一貫。問工墨錢幾何？\\

  \textbf{答曰：}十貫。\\

  \textbf{術曰：}置課鈔為實，以百貫為法，除之，即得。

}{%

  \textbf{Problem 96.}

  Suppose there are 1,000 guan of tax cash; the government charges 1 guan

  processing fee (工墨錢) per 100 guan. How much is the processing fee?\\

  \textbf{Answer:} 10 guan.\\

  \textbf{Method:} Place the tax cash as dividend; use 100 guan as divisor;

  divide; obtain.\\

  (Calculation: $1{,}000 \div 100 = 10$ guan processing fee)

}



\longline





%% ========================================================================

%% Chapter 8: 折變互差門

%% ========================================================================

\bilingualsubsection{折變互差門}{Conversion and Mutual Difference}

\label{sec:zhebian}



\bilingual{%

  \begin{rhymecn}

  折變互差要詳知，多少相乘作實推，\\

  以少求多因作母，除來便見兩無虧。

  \end{rhymecn}

  折變互差者，以物易物、以多變少之法也。

  其法以所求率與所有數相乘為實，以所知率為法除之，

  得所求之數。凡十五問。

}{%

  \begin{rhymecn}

  For conversion and mutual difference, you must know in detail;\\

  Multiply the quantities to make the dividend;\\

  To seek more from less, use multiplication as the base;\\

  Divide and see that nothing is lost.

  \end{rhymecn}

  This chapter covers problems of conversion between commodities,

  alloy calculations, price differentials, and weighted distribution.

  These problems involve multiple steps of multiplication and division,

  often with fractional rates.

}



\bilingual{%

  \textbf{問97.}

  今有香油三兩，折菜油四兩，只云香油斤價四百文。問菜油斤價幾何？\\

  \textbf{答曰：}如術計之。\\

  \textbf{術曰：}列香油數為實，以折率乘之，又以香油價為法，除之，即得。

}{%

  \textbf{Problem 97.}

  Suppose 3 liang of sesame oil is equivalent to 4 liang of vegetable oil;

  sesame oil costs 400 cash per jin. What is the price of vegetable oil per jin?\\

  \textbf{Answer:} Calculate according to the method.\\

  \textbf{Method:} Place the sesame oil quantity as dividend; multiply by the

  conversion rate; use sesame oil price as divisor; divide; obtain.

}



\bilingual{%

  \textbf{問98.}

  今有香油三兩，折菜油四兩，只云菜油八十四斤一十二兩，

  問香油幾何？\\

  \textbf{答曰：}如術計之。\\

  \textbf{術曰：}列菜油數為實，以折率除之，即得。

}{%

  \textbf{Problem 98.}

  Suppose 3 liang of sesame oil = 4 liang of vegetable oil;

  given 84 jin 12 liang of vegetable oil, how much sesame oil?\\

  \textbf{Answer:} Calculate according to the method.\\

  \textbf{Method:} Place the vegetable oil quantity as dividend; divide by the

  conversion rate; obtain.

}



\bilingual{%

  \textbf{問99.}

  今有麵四百單七斤，每六斤八分七釐半直錢一百文。問直錢幾何？\\

  \textbf{答曰：}如術計之。\\

  \textbf{術曰：}列麵數為實，以六斤八分七釐半為法，除之，即得。

}{%

  \textbf{Problem 99.}

  Suppose there are 407 jin of flour; 6 jin 8 fen 7.5 li costs 100 cash.

  How much cash?\\

  \textbf{Answer:} Calculate according to the method.\\

  \textbf{Method:} Place the flour quantity as dividend; use 6 jin 8 fen 7.5 li

  as divisor; divide; obtain.

}



\bilingual{%

  \textbf{問100.}

  今有油六碩八斗四升，每三斤七分五釐直錢一百文。問直錢幾何？\\

  \textbf{答曰：}如術計之。\\

  \textbf{術曰：}列油數為實，以三斤七分五釐為法，除之，即得。

}{%

  \textbf{Problem 100.}

  Suppose there are 6 shuo 8 dou 4 sheng of oil; 3 jin 7 fen 5 li costs 100 cash.

  How much cash?\\

  \textbf{Answer:} Calculate according to the method.\\

  \textbf{Method:} Place the oil quantity as dividend; use 3 jin 7 fen 5 li as

  divisor; divide; obtain.

}



\bilingual{%

  \textbf{問101.}

  今有麵數為實，以六斤八分七釐半除之。問得幾何？\\

  \textbf{術曰：}列麵數為實，以六斤八分七釐半為法，除之，即得。

}{%

  \textbf{Problem 101.}

  Suppose there is a quantity of flour as dividend; divide by 6 jin 8 fen 7.5 li.

  What is obtained?\\

  \textbf{Method:} Place the flour quantity as dividend; use 6 jin 8 fen 7.5 li as

  divisor; divide; obtain.

}



\bilingual{%

  \textbf{問102.}

  今有羅三十四匹六尺，每尺一百二十文。問計錢幾何？\\

  \textbf{答曰：}四貫一百五十二文。\\

  \textbf{術曰：}列羅數通為尺，以尺價乘之，即得。

}{%

  \textbf{Problem 102.}

  Suppose there are 34 bolts 6 chi of gauze, each chi costing 120 cash.

  How much cash in total?\\

  \textbf{Answer:} 4 guan 152 cash.\\

  \textbf{Method:} Convert the gauze quantity entirely to chi; multiply by chi-price;

  obtain.

}



\bilingual{%

  \textbf{問103.}

  今有絹七匹，每匹四貫。問計錢幾何？\\

  \textbf{答曰：}二十八貫。\\

  \textbf{術曰：}列絹數，以匹價乘之，即得。

}{%

  \textbf{Problem 103.}

  Suppose there are 7 bolts of silk, each costing 4 guan.

  How much cash in total?\\

  \textbf{Answer:} 28 guan.\\

  \textbf{Method:} Place the silk quantity; multiply by bolt-price; obtain.\\

  (Calculation: $7 \times 4 = 28$)

}



\bilingual{%

  \textbf{問104.}

  今有絲一斤價錢二貫文，問一兩價錢幾何？\\

  \textbf{答曰：}一百二十五文。\\

  \textbf{術曰：}置絲斤價為實，以十六兩為法，除之，即得。

}{%

  \textbf{Problem 104.}

  Suppose silk costs 2 guan per jin. What is the price per liang?\\

  \textbf{Answer:} 125 cash per liang.\\

  \textbf{Method:} Place the silk jin-price as dividend; use 16 liang as divisor;

  divide; obtain.\\

  (Calculation: $2{,}000 \div 16 = 125$)

}



\bilingual{%

  \textbf{問105.}

  今有田一畝，闊十五步，長十六步。問為頃畝幾何？\\

  \textbf{答曰：}一畝。\\

  \textbf{術曰：}列闊步、長步相乘得積步，以畝法二百四十步除之，即得。

}{%

  \textbf{Problem 105.}

  Suppose there is a field of 1 mu, 15 paces wide and 16 paces long.

  What is it in qing and mu?\\

  \textbf{Answer:} 1 mu.\\

  \textbf{Method:} Place width and length in paces; multiply to get accumulated paces;

  divide by the mu-conversion of 240 paces; obtain.\\

  (Calculation: $15 \times 16 = 240$ paces; $240 \div 240 = 1$ mu)

}



\bilingual{%

  \textbf{問106.}

  今有粟九斗，每斛直錢一貫二百文。問直錢幾何？\\

  \textbf{答曰：}一貫八十文。\\

  \textbf{術曰：}置粟數為實，以斛價為法，乘之，即得。

}{%

  \textbf{Problem 106.}

  Suppose there are 9 dou of millet, each hu costing 1 guan 200 cash.

  What is the value?\\

  \textbf{Answer:} 1 guan 80 cash.\\

  \textbf{Method:} Place the millet quantity as dividend; use hu-price as divisor;

  multiply; obtain.\\

  (9 dou = 0.9 hu; $0.9 \times 1{,}200 = 1{,}080$)

}



\bilingual{%

  \textbf{問107.}

  今有金五十兩，令二八共造，問各幾何？\\

  \textbf{答曰：}金四十兩，銀一十兩。\\

  \textbf{術曰：}置共金為實，以分母為法，除之，各以分子乘之，即得。

}{%

  \textbf{Problem 107.}

  Suppose there are 50 liang of metal, to be made in a 2:8 ratio.

  How much of each?\\

  \textbf{Answer:} Gold 40 liang, silver 10 liang.\\

  \textbf{Method:} Place the total metal as dividend; use the denominator as divisor;

  divide; multiply each by its numerator; obtain.\\

  (``二八共造'' = 2 parts silver, 8 parts gold; $50 \times \frac{8}{10} = 40$ gold,

  $50 \times \frac{2}{10} = 10$ silver)

}



\bilingual{%

  \textbf{問108.}

  今有銀二十四兩，直錢三貫六百文。只有銀五十九兩八銖。

  問直錢幾何？\\

  \textbf{答曰：}八十九貫七百文。\\

  \textbf{術曰：}列只有銀數，以乘原直錢為實，以原銀為法，除之，即得。

}{%

  \textbf{Problem 108.}

  Suppose 24 liang of silver is worth 3 guan 600 cash.

  If one has 59 liang 8 zhu of silver, what is it worth?\\

  \textbf{Answer:} 89 guan 700 cash.\\

  \textbf{Method:} Place the available silver; multiply by original cash value as

  dividend; use original silver as divisor; divide; obtain.

}



\bilingual{%

  \textbf{問109.}

  今有絲八斤二兩二十一銖，織錦七匹六尺五寸。只有絲九十五斤三兩六銖。

  問織錦幾何？\\

  \textbf{術曰：}（匹法同前）列只有絲數，以乘錦數為實，以原絲為法，除之，即得。

}{%

  \textbf{Problem 109.}

  Suppose 8 jin 2 liang 21 zhu of silk weaves 7 bolts 6 chi 5 cun of brocade.

  If one has 95 jin 3 liang 6 zhu of silk, how much brocade can be woven?\\

  \textbf{Method:} (Bolt conversion same as before) Place the available silk;

  multiply by brocade quantity as dividend; use original silk as divisor;

  divide; obtain.

}



\bilingual{%

  \textbf{問110.}

  今有錢五貫八百三十文，買麩八斗五升。只有錢三十七貫七百六十文。

  問買麩幾何？\\

  \textbf{答曰：}五碩五斗二升。\\

  \textbf{術曰：}列只有錢數，以乘原麩為實，以原錢為法，除之，即得。

}{%

  \textbf{Problem 110.}

  Suppose 5 guan 830 cash buys 8 dou 5 sheng of bran.

  If one has 37 guan 760 cash, how much bran can be bought?\\

  \textbf{Answer:} 5 shuo 5 dou 2 sheng.\\

  \textbf{Method:} Place the available cash; multiply by original bran as dividend;

  use original cash as divisor; divide; obtain.

}



\bilingual{%

  \textbf{問111.}

  今有錢一貫一百七十五文，買麥五斗四升。只有錢二十三貫五百文。

  問買麥幾何？\\

  \textbf{答曰：}十碩八斗。\\

  \textbf{術曰：}列只有錢數，以乘原麥為實，以原錢為法，除之，即得。

}{%

  \textbf{Problem 111.}

  Suppose 1 guan 175 cash buys 5 dou 4 sheng of wheat.

  If one has 23 guan 500 cash, how much wheat can be bought?\\

  \textbf{Answer:} 10 shuo 8 dou.\\

  \textbf{Method:} Place available cash; multiply by original wheat as dividend;

  use original cash as divisor; divide; obtain.

}



\longline



%% ========================================================================

%% APPENDIX: Textual Notes

%% ========================================================================

\newpage

\bilingualsection{Appendix: Textual Notes}{附錄：技術說明}



\bilingual{%

  \textbf{一、算籌記數法}



  中國古代以算籌記數，縱橫相間，以別位值。

  一縱十橫，百立千僵。算籌以竹為之，長約三寸，徑約一分。

  記數之法：一至四以縱置之，五以一橫置之，六至九以橫五加縱數置之。

  十位則一橫十縱，百位復為縱，千位復為橫，以此相間。

  此訣所謂「一縱十橫，百立千僵」也。



  \smallskip

  \textbf{二、增乘開方法}



  增乘開方者，以漸進之法求方程之根也。

  其法置積為實，立初商，以次乘實，遞減求之。

  賈憲首創此術，朱世傑廣而用之，以解高次方程。

  此法與西洋之魯菲尼—霍納法相合，而早於西洋數百年。



  \smallskip

  \textbf{三、天元術}



  天元術者，以天元一為未知數，立方程而求解之術也。

  其法以「太」字標常數項，以元之冪次遞降排列。

  李冶首創於《測圓海鏡》，朱世傑繼之於《四元玉鑑》，

  並推廣為四元術。此為中國代數學之最高成就。



  \smallskip

  \textbf{四、本書流傳}



  算學啓蒙成書於元大德三年（1299年），在中國久已失傳。

  清初始從朝鮮發現，道光十九年（1839年）揚州羅士琳重刻之。

  此書東傳日本，影響和算之發展甚大。

}{%

  \textbf{I. Counting Rod Notation (算籌記數法)}



  Ancient Chinese mathematics used counting rods (算籌) made of bamboo,

  approximately 3 inches long, arranged on a counting board. The notation

  alternates between vertical and horizontal orientations by place value:

  units vertical, tens horizontal, hundreds vertical, thousands horizontal.

  Digits 1--4 are represented by 1--4 rods in the appropriate orientation;

  5 is a single horizontal bar (for units) or vertical bar (for tens);

  6--9 combine the 5-bar with 1--4 additional rods. This is the system

  described by the rhyme ``units vertical, tens horizontal'' (一縱十橫).



  \smallskip

  \textbf{II. Incremental Root Extraction (增乘開方法)}



  The ``incremental multiplication'' method for root extraction was developed

  by Jia Xian (11th century) and extended by Zhu Shijie. It systematically

  extracts roots digit by digit through a process of successive approximation.

  This method is algebraically equivalent to the Ruffini--Horner method

  developed in Europe centuries later.



  \smallskip

  \textbf{III. The Tian Yuan Method (天元術)}



  The ``celestial element'' method uses tian yuan yi (天元一, ``celestial element one'')

  as the unknown variable. Equations are set up with coefficients arranged in

  ascending or descending powers. Li Ye (李冶) pioneered this in his

  \textit{Ce Yuan Hai Jing} (測圓海鏡, 1248); Zhu Shijie extended it to

  the si yuan (四元, ``four elements'') method for systems of equations

  in his \textit{Siyuan Yujian} (四元玉鑑, 1303).



  \smallskip

  \textbf{IV. Transmission of the Text}



  The \textit{Suanxue Qimeng} was written in 1299 and was lost in China for centuries.

  A Korean edition was rediscovered in the early Qing dynasty. Luo Shilin (羅士琳)

  republished it in Yangzhou in 1839. The work profoundly influenced the

  development of \textit{wasan} (和算, Japanese mathematics).

}



\vfill

\begin{center}

\textcolor{rulecolor}{\rule{0.3\textwidth}{0.4pt}}\\[4pt]

{\small\textit{End of Volume I (卷上) --- 113 Problems}}\\[2pt]

{\small Bilingual Edition --- \LaTeX\ Modern LaTeX edition}

\end{center}



\end{document}

